\documentclass[12pt, a4paper, reqno, oneside]{book}
\usepackage[left=2.5cm, right=2cm, top=2cm, bottom=2cm]{geometry}
\usepackage{amsmath,amsxtra,amssymb,latexsym,amscd,amsthm}
\usepackage{microtype}
\usepackage{indentfirst}
\usepackage{color}
\usepackage{empheq}
\usepackage[USenglish]{babel}
\usepackage[mathscr]{eucal}
\usepackage{amsfonts}
\usepackage{graphicx}
\usepackage{float}
\usepackage{listings}
\usepackage[dvipsnames]{xcolor}
\usepackage{verbatim}
\usepackage{tikz}
\usetikzlibrary{calc,positioning,shapes.geometric,fit,arrows.meta,shadows.blur}
\usepackage{array}
\usepackage{multirow}
\usepackage{algorithm}
\usepackage{algpseudocode}
\usepackage{subfigure}
\usepackage{pgfplots}
\pgfplotsset{compat=1.18}
\usepackage{titlesec}
\usepackage{fancyhdr}
\usepackage{rotating}
\usepackage{longtable}
\usepackage{supertabular}
\fancyhf{}
\fancyhead[LO,RE]{\small{\leftmark}}
\fancyfoot[C]{\thepage}
\usepackage{ulem}
\usepackage{pdfpages}
\usepackage[square,numbers]{natbib}
\usepackage{subcaption,booktabs}
\usepackage{chngcntr}
\usepackage{tabularray}
\usepackage{tabularx}
\usepackage[labelfont=bf, textfont=it, justification=centering]{caption}
\captionsetup{labelfont={color=blue, bf}, textfont={color=black, it}}
\renewcommand{\labelitemi}{+}
\hyphenation{op-tical net-works semi-conduc-tor}
\usepackage[hyphens,spaces,obeyspaces]{url}
\usepackage[colorlinks,bookmarksopen,bookmarksnumbered,
            citecolor=blue,colorlinks=true,
            linkcolor=blue,urlcolor=blue]{hyperref}

\renewcommand{\baselinestretch}{1.25}

%─── unnumbered chapter/section helpers (matches template) ─────────────
\newcommand\chap[1]{%
  \chapter*{#1}%
  \addcontentsline{toc}{chapter}{#1}%
  \markboth{#1}{}}
\newcommand\sect[1]{%
  \section*{#1}%
  \addcontentsline{toc}{section}{#1}}

%─── counter formatting ────────────────────────────────────────────────
\counterwithin{algorithm}{chapter}
\renewcommand{\thealgorithm}{\thechapter.\arabic{algorithm}}
\renewcommand{\theequation}{\arabic{chapter}.\arabic{equation}}

%─── theorem environments ─────────────────────────────────────────────
\newtheorem{lemma}{\bf Lemma}[chapter]
\newtheorem{theorem}{\bf Theorem}[chapter]
\newtheorem{proposition}{\bf Proposition}[chapter]
\newtheorem{corollary}{\bf Corollary}[chapter]
\newtheorem{definition}{\bf Definition}[chapter]
\newtheorem{remark}{\bf Remark}[section]
\newtheorem{assumption}{\bf Assumption}[chapter]
\newtheorem{invariant}{\bf Invariant}[chapter]

%─── project-specific macros ──────────────────────────────────────────
\newcommand{\padonap}{\textsc{PADOM}}
\newcommand{\tier}[1]{$T_{#1}$}

%─── subsubsubsection support ─────────────────────────────────────────
\titleclass{\subsubsubsection}{straight}[\subsection]
\newcounter{subsubsubsection}[subsubsection]
\renewcommand\thesubsubsubsection{\thesubsubsection.\arabic{subsubsubsection}}
\renewcommand\theparagraph{\thesubsubsubsection.\arabic{paragraph}}
\titleformat{\subsubsubsection}
  {\normalfont\normalsize\bfseries}{\thesubsubsubsection}{1em}{}
\titlespacing*{\subsubsubsection}{0pt}{3.25ex plus 1ex minus .2ex}{1.5ex plus .2ex}
\makeatletter
\renewcommand\paragraph{\@startsection{paragraph}{5}{\z@}%
  {3.25ex \@plus1ex \@minus.2ex}{-1em}{\normalfont\normalsize\bfseries}}
\renewcommand\subparagraph{\@startsection{subparagraph}{6}{\parindent}%
  {3.25ex \@plus1ex \@minus .2ex}{-1em}{\normalfont\normalsize\bfseries}}
\def\toclevel@subsubsubsection{4}
\def\toclevel@paragraph{5}
\def\l@subsubsubsection{\@dottedtocline{4}{7em}{4em}}
\def\l@paragraph{\@dottedtocline{5}{10em}{5em}}
\def\l@subparagraph{\@dottedtocline{6}{14em}{6em}}
\makeatother
\setcounter{secnumdepth}{4}
\setcounter{tocdepth}{4}

%─── graphics path ────────────────────────────────────────────────────
\graphicspath{{./figs/}}

%─── listing styles ───────────────────────────────────────────────────
\definecolor{mygreen}{rgb}{0,0.6,0}
\definecolor{mygray}{rgb}{0.5,0.5,0.5}
\definecolor{mymauve}{rgb}{0.58,0,0.82}

\lstdefinestyle{custompython}{
  belowcaptionskip=1\baselineskip,
  breaklines=true, frame=single,
  xleftmargin=\parindent,
  language=Python,
  showstringspaces=false,
  basicstyle=\footnotesize\ttfamily,
  keywordstyle=\bfseries\color{green!40!black},
  commentstyle=\itshape\color{purple!40!black},
  identifierstyle=\color{blue},
  stringstyle=\color{orange},
  numbers=left, numberstyle=\tiny\color{gray}, stepnumber=1,
}

\definecolor{eclipseStrings}{RGB}{42,0,255}
\definecolor{eclipseKeywords}{RGB}{127,0,85}
\colorlet{numb}{magenta!60!black}
\lstdefinelanguage{json}{
  basicstyle=\normalfont\ttfamily,
  commentstyle=\color{eclipseStrings},
  stringstyle=\color{eclipseKeywords},
  numbers=left, numberstyle=\scriptsize, stepnumber=1, numbersep=8pt,
  showstringspaces=false, breaklines=true, frame=lines,
  string=[s]{"}{"},
  comment=[l]{:\ "}, morecomment=[l]{:"},
  literate=*{0}{{\color{numb}0}}{1}{1}{{\color{numb}1}}{1}
            {2}{{\color{numb}2}}{1}{3}{{\color{numb}3}}{1}
            {4}{{\color{numb}4}}{1}{5}{{\color{numb}5}}{1}
            {6}{{\color{numb}6}}{1}{7}{{\color{numb}7}}{1}
            {8}{{\color{numb}8}}{1}{9}{{\color{numb}9}}{1},
}
\lstset{style=custompython}

\newcolumntype{C}[1]{>{\centering\arraybackslash}p{#1}}
\definecolor{processblue}{cmyk}{0.96,0,0,0}

%─────────────────────────────────────────────────────────────────────
\begin{document}
\hypersetup{pageanchor=false}

%══════════════════════════════════════════════════════════════════════
%  TITLE PAGE 1
%══════════════════════════════════════════════════════════════════════
\begin{titlepage}
\begin{center}
  \begin{tikzpicture}[overlay,remember picture]
    \draw[line width=3pt,rounded corners=0pt]
      ($(current page.north west)+(20mm,-20mm)$)
      rectangle
      ($(current page.south east)+(-15mm,20mm)$);
    \draw[line width=1pt,rounded corners=0pt]
      ($(current page.north west)+(21.5mm,-21.5mm)$)
      rectangle
      ($(current page.south east)+(-16.5mm,21.5mm)$);
  \end{tikzpicture}
  \\[1mm]
  \centerline{\bf VIETNAM NATIONAL UNIVERSITY HANOI}
  \centerline{\bf UNIVERSITY OF ENGINEERING AND TECHNOLOGY}
  \vspace*{1.8cm}
  \begin{figure}[H]
    \centering
    \IfFileExists{figs/uet.jpg}{%
      \includegraphics[width=0.23\linewidth]{figs/uet.jpg}%
    }{%
      \fbox{\parbox[c][2.2cm][c]{0.23\linewidth}{\centering UET Logo}}%
    }
  \end{figure}
  \vspace*{1cm}
  \centerline{\bf KHONG MANH TUAN}
  \vspace*{2.5cm}
  \centerline{\bf AI-AUGMENTED NFV ORCHESTRATION FOR}
  \centerline{\bf PROACTIVE DDOS MITIGATION IN DATA CENTER NETWORKS}
  \vspace*{2.5cm}
  \centerline{\bf GRADUATION THESIS}
  \vspace*{0.2cm}
  \centerline{\bf Major: Computer Science}
  \vfill
  \centerline{\bf Hanoi -- 2026}
  \vspace{13pt}
\end{center}
\end{titlepage}

\newpage

\begin{titlepage}
\begin{center}
  \begin{tikzpicture}[overlay,remember picture]
    \draw[line width=3pt,rounded corners=0pt]
      ($(current page.north west)+(20mm,-20mm)$)
      rectangle
      ($(current page.south east)+(-15mm,20mm)$);
    \draw[line width=1pt,rounded corners=0pt]
      ($(current page.north west)+(21.5mm,-21.5mm)$)
      rectangle
      ($(current page.south east)+(-16.5mm,21.5mm)$);
  \end{tikzpicture}
  \\[1mm]
  \centerline{\bf VIETNAM NATIONAL UNIVERSITY HANOI}
  \centerline{\bf UNIVERSITY OF ENGINEERING AND TECHNOLOGY}
  \vspace*{1.8cm}
  \begin{figure}[H]
    \centering
    \IfFileExists{figs/uet.jpg}{%
      \includegraphics[width=0.23\linewidth]{figs/uet.jpg}%
    }{%
      \fbox{\parbox[c][2.2cm][c]{0.23\linewidth}{\centering UET Logo}}%
    }
  \end{figure}
  \vspace*{0.9cm}
  \centerline{\bf KHONG MANH TUAN}
  \vspace*{2.2cm}
  \centerline{\bf AI-AUGMENTED NFV ORCHESTRATION FOR}
  \centerline{\bf PROACTIVE DDOS MITIGATION IN DATA CENTER NETWORKS}
  \vspace*{1.7cm}
  \centerline{\bf Major: COMPUTER SCIENCE}
  \vspace*{1.6cm}
  \centerline{\bf Supervisor:}
  \centerline{\bf Assoc. Prof. Nguyen Ngoc Hoa}
  \centerline{\bf PhD. Du Phuong Hanh}
  \vfill
  \centerline{\bf Hanoi -- 2026}
  \vspace{13pt}
\end{center}
\end{titlepage}

\newpage
\pagenumbering{gobble}


\chapter*{Acknowledgements}
\vspace{20pt}

I would like to express my sincere gratitude to my supervisor for the invaluable guidance, feedback, and encouragement provided throughout the course of this research. Their expertise in network security and AI-driven systems was instrumental in shaping both the direction and quality of this thesis.

I am also grateful to the faculty members of the Department of Computer Networks and Data Communication at the University of Engineering and Technology for their support and for the strong academic foundation they provided.

I acknowledge that this work may still have room for improvement, and I welcome any constructive feedback that could further strengthen the research in future iterations.

Sincere gratitude!

\newpage

\chapter*{Authorship}
\vspace{20pt}

I declare that this thesis is my original work, conducted under the supervision of my academic advisor at the University of Engineering and Technology, Vietnam National University Hanoi. This work has not been submitted for any other degree or professional qualification.

All references and related works used in this thesis have been appropriately acknowledged and cited. The quantitative results reported in this thesis include measured AI and Mininet observations as well as explicitly marked stub-simulated or estimated orchestration measurements from the \padonap{} prototype.

I affirm that this declaration is made with full integrity and in accordance with the institution's academic policies.

\begin{flushright}
{\it Author}\\[2cm]
{\bf Khong Manh Tuan}
\end{flushright}
\newpage

\chapter*{Supervisor's Approval}
\vspace{20pt}

I hereby approve that the thesis in its current form is ready for committee examination as a requirement for the Engineer of Computer Networks and Data Communication degree at the VNU University of Engineering and Technology.

All references from relevant studies are explicitly sourced from the list of references in the thesis. The thesis does not copy documents or research works of others without specifying the references.

\begin{flushright}
{\it Signature}\\[2cm]
{\bf Nguyen Ngoc Hoa}
{\bf }
\end{flushright}

\newpage

\chapter*{Abstract}
\addcontentsline{toc}{chapter}{Abstract}
\vspace{20pt}

Distributed denial-of-service (DDoS) attacks pose a serious threat to data center networks because they can rapidly deplete link capacity, transport layer state, and shared mitigation resources. However, the most modern closed-loop defense mechanisms in network functional virtualization (NFV) environments remain largely reactive, exhibiting end-to-end mitigation latency that fails to keep pace with the speed of modern large-scale attacks. To overcome this gap, this thesis proposes \padonap{}, a four-stage closed-loop architecture inspired by ONAP.

Unlike conventional methods, \padonap{} incorporates AI strategies in detection and prediction. A lightweight \textbf{XGBoost} classifier on a 17-dimensional window-level feature vector provides explainable, reactive detection on the six-class CICDDoS2019 label space used in the Track~A experiment. Simultaneously, a \textbf{CNN + Bidirectional LSTM with attention pooling} (CNN+BiLSTM+Attention) acts as a multi-horizon predictor at 60, 300, and 600 seconds to proactively anticipate attack risk; this architecture was chosen after benchmarking five sequence-modelling families against a persistence baseline and a tabular XGBoost-flat baseline on CICIDS2017. The outputs of this AI are mapped to a 5-tier hierarchical orchestration strategy within the ONAP-inspired Policy Framework, enabling proactive pre-positioning and dynamic scaling of Virtual Network Functions (VNFs) for data filtering. Furthermore, the system integrates SHapley Additive exPlanations (SHAP) to provide explainable policy metadata, effectively addressing the "black box" limitations of traditional AI defenses.

Evaluated on the CICDDoS2019 benchmark dataset, \padonap{} achieves robust classification performance while trading maximum theoretical accuracy for practical NFV deployment capabilities. The results show that \padonap{} minimizes end-to-end latency compared to existing basic response methods, supports SLA-compliant allocation in the evaluated capacity setups, and provides automated, transparent security orchestration for the data center network infrastructure.

\newpage

\chapter*{List of Abbreviations}
\addcontentsline{toc}{chapter}{List of Abbreviations}
\vspace{10pt}

\begin{table}[H]
\centering
\begin{tabular}{|m{3cm}|m{9cm}|}
\hline
\textbf{Abbreviation} & \textbf{Full Form} \\ \hline
AI    & Artificial Intelligence \\ \hline
AUC   & Area Under the ROC Curve  \\ \hline
DCN   & Data Center Network \\ \hline
DDoS  & Distributed Denial-of-Service  \\ \hline
DL    & Deep Learning \\ \hline
gNMI  & gRPC Network Management Interface \\ \hline
LP    & Linear Program / Linear Programming \\ \hline
LSTM  & Long Short-Term Memory \\ \hline
MANO  & Management and Orchestration \\ \hline
ML    & Machine Learning \\ \hline
NFV   & Network Function Virtualization \\ \hline
ONAP  & Open Network Automation Platform \\ \hline
SFC   & Service Function Chaining \\ \hline
SHAP  & SHapley Additive exPlanations \\ \hline
SLA   & Service Level Agreement \\ \hline
SO    & Service Orchestrator (ONAP component) \\ \hline
VNF   & Virtual Network Function \\ \hline
XGB   & XGBoost (eXtreme Gradient Boosting) \\ \hline
BiLSTM & Bidirectional Long Short-Term Memory \\ \hline
CNF & Cloud-Native Network Function \\ \hline
CNN & Convolutional Neural Network \\ \hline
DCAE & Data Collection, Analytics, and Events \\ \hline
GRU & Gated Recurrent Unit \\ \hline
HTTP & Hypertext Transfer Protocol \\ \hline
ICMP & Internet Control Message Protocol \\ \hline
MLP & Multi-Layer Perceptron \\ \hline
OOM & ONAP Operations Manager \\ \hline
OVS & Open vSwitch \\ \hline
PAP & Policy Administration Point \\ \hline
PCAP & Packet Capture \\ \hline
PDP & Policy Decision Point \\ \hline
RPC & Remote Procedure Call \\ \hline
\end{tabular}
\end{table}

\newpage

\tableofcontents
\newpage
\listoffigures
\newpage
\listoftables
\begingroup
\listofalgorithms
\let\clearpage\relax
\endgroup

\newpage
\pagenumbering{arabic}

\chapter{Introduction}
\label{ch:introduction}

\section{Motivation}
Modern data center networks host latency-sensitive and mission-critical
services that require continuous availability. Distributed Denial-of-Service
(DDoS) attacks remain a persistent threat because they can exhaust link
capacity, transport-layer state, and application-layer resources through
coordinated traffic sources and heterogeneous attack vectors
\cite{sharafaldin2019cicddos}. In multi-tenant data center environments, this
threat is amplified by resource sharing: an attack against one service
may consume shared compute, network, and security-function capacity, thereby
degrading legitimate tenants that are not direct attack targets.

Network Function Virtualization (NFV) and Service Function Chaining (SFC)
provide mechanisms for deploying and steering traffic through software-based
security functions, such as firewalls, rate limiters, and traffic scrubbers
\cite{etsi_nfv_architecture,sfc_rfc7665}. Cloud-native deployment platforms
further simplify the packaging, scaling, and observability of these functions
\cite{kubernetes_docs}. However, elasticity alone does not eliminate the
attack-to-enforcement latency of a reactive defense loop. Such a loop must
observe suspicious traffic, classify or score the event, trigger a policy
decision, submit an orchestration request, wait until the mitigation function
is ready, and finally enforce traffic steering or filtering. During this
exposure window, attack traffic may continue to degrade service quality,
consume shared resources, and increase the probability of SLA violation.

The motivation of this thesis lies in the operational gap between AI-based DDoS analytics and orchestrated mitigation. Many learning-based approaches report classification accuracy, F1-score, or forecasting error in isolation. However, the key operational question is whether an early warning signal provides enough lead time for mitigation, whether the decision is explainable to policy logic, and whether the resulting closed loop reduces service degradation under orchestration delay. To address this gap, PADOM combines near-real-time flow-level detection, multi-horizon forecasting, explainable policy metadata, and an ONAP-inspired closed loop. \cite{onap_docs}.


\section{Research Challenges}

The development and evaluation of \padonap{} address several key research
challenges:

\begin{itemize}
    \item \textbf{Q1: Proactive Mitigation under Orchestration Delay.}
    How can a DDoS defense system reduce mitigation delay in NFV/CNF environments? Reactive systems usually act only after attack evidence exceeds a predefined threshold. In practice, the response loop also includes telemetry collection, detection, policy evaluation, orchestration, CNF readiness, and traffic enforcement. This research examines whether early warning signals can prepare mitigation functions before the attack reaches its most damaging stage.

    \item \textbf{Q2: Accurate Forecasting with Controlled False Positives.}
    How can multi-horizon forecasting provide proactive defense value without
    causing unnecessary mitigation actions during benign traffic fluctuations?
    Forecasting can create lead time for mitigation, but incorrect escalation may
    waste compute and network resources or disturb legitimate traffic. The
    challenge is to combine near-real-time DDoS detection with short-horizon
    forecasting and confidence-aware escalation rules, so that proactive actions
    are triggered only when the predicted risk is sufficiently reliable.

    \item \textbf{Q3: Explainable Translation from AI Outputs to Policy Actions.}
    How can AI detection and forecasting outputs be converted into transparent
    policy decisions that can be executed by the orchestration layer? A practical
    defense system must not only produce attack probabilities or predicted
    traffic risks, but also map them to operational response tiers such as
    alerting, preemption, mitigation, and isolation. This research must therefore
    define a policy mechanism that records why a response tier was selected and
    exposes interpretable metadata for auditability.

    \item \textbf{Q4: ONAP-Compatible Closed-Loop CNF Orchestration.}
    How can a lightweight ONAP-style closed loop connect telemetry analytics,
    policy decisions, service orchestration, and CNF enforcement without claiming
    a full ONAP deployment? \padonap{} must align with the relevant ONAP
    closed-loop roles while clearly distinguishing implemented components from
    simulated or proposed components.

    \item \textbf{Q5: Evidence-Based Evaluation of Mitigation and SLA Impact.}
    How can the system be evaluated without overstating unmeasured deployment
    results? The core evaluation should measure AI detection performance,
    forecasting behavior, forecasting lead time, closed-loop decision latency,
    CNF readiness time, and mitigation effectiveness. Extended evaluation should
    consider resource overhead and tenant-level SLA impact where the testbed
    supports such measurements.
\end{itemize}

\section{Objective and Contents}

The primary objective of this thesis is to design, implement, and evaluate
\padonap{}, an AI-augmented NFV/CNF orchestration framework for proactive
DDoS mitigation in data center networks. The thesis investigates how
near-real-time detection, multi-horizon attack forecasting, explainable policy
decisions, and CNF-based service orchestration can be combined to reduce the
attack-to-enforcement latency of traditional reactive mitigation systems.

Specifically, \padonap{} aims to shift the defense paradigm from purely reactive
mitigation to proactive mitigation preparation. Instead of waiting until attack
traffic fully exceeds a mitigation threshold, the proposed framework uses both
immediate flow-level detection and short-horizon traffic forecasting to generate
early warning signals. These signals are then translated into graduated policy
actions, such as alerting, pre-positioning, mitigation, and isolation, through an
ONAP-inspired closed-loop architecture. The framework further considers
multi-tenant service requirements by incorporating resource-aware decision logic
for mitigation deployment and traffic handling.

To achieve this objective, the thesis covers the following contents:

\begin{itemize}
    \item A detailed analysis of DDoS mitigation challenges in data center
    networks, with emphasis on multi-tenant resource sharing, orchestration delay,
    CNF/VNF readiness time, and the limitations of purely reactive defense loops.

    \item The design of the \padonap{} architecture, including telemetry
    collection, AI-based detection and forecasting, explainable policy decision
    logic, closed-loop orchestration, and CNF enforcement.

    \item The development of a dual-track AI analytics engine that combines an
    immediate multi-class DDoS detection model with a multi-horizon forecasting
    model. The detection component identifies current attack states, while the
    forecasting component estimates near-future traffic risk to support proactive
    mitigation preparation.

    \item The design of a policy translation mechanism that maps AI outputs into
    auditable response tiers. This mechanism records relevant decision metadata,
    such as detection confidence, forecasted risk, selected mitigation tier, and
    policy rationale, to improve transparency and operational traceability.

    \item The implementation of a lightweight orchestration loop that connects
    telemetry ingestion, policy decision making, orchestration control, and
    CNF-based mitigation enforcement in a Kubernetes environment.

\end{itemize}

\section{Contribution Highlights}

This thesis provides the following key contributions to the field of network
security and NFV/CNF orchestration:

\begin{enumerate}
    \item \textbf{Proactive DDoS Mitigation through Multi-Horizon Forecasting.}
    PADOM integrates near-real-time DDoS detection with multi-horizon
    forecasting to generate early warning signals before attack traffic reaches
    its most damaging stage. These forecasts support proactive mitigation
    preparation, including alerting, pre-positioning, and pre-scaling lightweight
    mitigation functions.

    \item \textbf{ONAP-Inspired Closed-Loop Orchestration for CNF-Based Defense.}
    The proposed framework implements a lightweight ONAP-inspired closed loop
    that links telemetry collection, AI analytics, policy decision making,
    orchestration control, and CNF enforcement. It also explicitly distinguishes
    implemented components from simulated or proposed ONAP-related functions.

    \item \textbf{Explainable Translation from AI Outputs to Mitigation Tiers.}
    PADOM defines a policy translation mechanism that maps detection confidence,
    forecasted risk, and traffic indicators to graduated response tiers, including
    alerting, pre-positioning, mitigation, and isolation. The mechanism records
    decision metadata to improve auditability, operational traceability, and
    post-event analysis.
\end{enumerate}

\section{Thesis Outline}

The remainder of this thesis is structured as follows:

\begin{itemize}
    \item \textbf{Chapter~\ref{ch:background}} introduces the fundamental
    background concepts required for this work, including DDoS attack taxonomy,
    data center network characteristics, NFV and SFC mechanisms, ONAP
    closed-loop orchestration concepts, Kubernetes-based CNF deployment, and
    machine learning principles for detection and forecasting.

    \item \textbf{Chapter~\ref{ch:related}} reviews related work on ML-based
    DDoS detection, traffic forecasting, NFV/SFC-based mitigation, and
    closed-loop security orchestration. The chapter identifies the research gap
    between AI-based attack analytics and operational mitigation under
    orchestration delay.

    \item \textbf{Chapter~\ref{ch:proposed}} presents the proposed \padonap{}
    framework. It describes the overall system architecture, telemetry and
    analytics pipeline, dual-track AI engine, explainable policy translation
    mechanism, resource-aware mitigation logic, and closed-loop orchestration
    workflow.

    \item \textbf{Chapter~\ref{ch:eval}} presents the experimental methodology
    and evaluation results. The chapter evaluates detection performance,
    forecasting behavior, proactive lead time, closed-loop latency, CNF readiness
    time, mitigation effectiveness, and resource overhead where supported by the
    testbed. It also discusses experimental limitations and distinguishes
    measured results from simulated or future-work components.
\end{itemize}

%══════════════════════════════════════════════════════════════════════
%  CHAPTER 2 — Background
%══════════════════════════════════════════════════════════════════════
%======================================================================
%======================================================================
\chapter{Background}
\label{ch:background}
%======================================================================

This chapter presents the foundational concepts required to understand the
design and evaluation of \padonap{}. The proposed framework combines
AI-based DDoS analytics, multi-horizon traffic forecasting, closed-loop
orchestration, and CNF enforcement. Therefore, this chapter first introduces
DDoS attacks in Data Center Networks (DCNs) and
explains why shared-resource environments increase the impact of such attacks.
It then describes the CICDDoS2019 dataset and the window-based traffic
representation used by the AI pipeline. Next, the chapter reviews SDN, NFV,
CNFs, and Service Function Chaining as the networking foundations for dynamic
mitigation. It then introduces ONAP closed-loop automation and explains how
\padonap{} adopts ONAP concepts without claiming a full ONAP deployment.
Finally, the chapter discusses XGBoost-based detection, CNN+BiLSTM+Attention
multi-horizon forecasting, SHAP explainability, and SLA-aware resource
allocation.

%----------------------------------------------------------------------
\section{DDoS Attacks in Data Center Networks}
\label{sec:ddos-bg}
%----------------------------------------------------------------------

Data Center Networks (DCNs) host latency-sensitive and mission-critical
services, including web applications, APIs, storage systems, enterprise
platforms, and internal network services. These services are expected to remain
available even when the network experiences sudden traffic surges or malicious
activity. A Distributed Denial-of-Service (DDoS) attack attempts to reduce
service availability by sending traffic from multiple sources toward a victim
service or infrastructure component. Such attacks may exhaust link bandwidth,
server processing capacity, connection state, queueing resources, or
application-layer resources.

In traditional isolated deployments, the impact of a DDoS attack is often
discussed mainly in terms of the direct victim. In DCNs, however, the effect
can spread beyond the attacked service because multiple applications, tenants,
or service domains may share the same physical and virtual infrastructure.
Compute nodes, switching capacity, ingress gateways, load balancers,
monitoring systems, and security functions may all be shared. Therefore, an
attack against one service domain can consume common resources and degrade
other legitimate services that are not direct attack targets. This
shared-resource coupling is one of the main reasons why DDoS mitigation in
DCNs must be studied not only as an attack classification problem, but also as
an orchestration and resource management problem.

% %---------------------- FIGURE PLACEHOLDER ----------------------------
% \begin{figure}[htbp]
%     \centering
%     % \includegraphics[width=0.95\linewidth]{figures/ddos_shared_resource_dcn.pdf}
%     \fbox{\parbox{0.9\linewidth}{
%     \centering
%     Placeholder for Figure: DDoS attack in a shared-resource data center
%     network. Show botnets sending traffic through the Internet/edge gateway
%     toward one protected service, while other services are affected indirectly
%     through shared bandwidth, compute, and security-function resources.
%     }}
%     \caption{DDoS impact in a shared-resource data center network. An attack
%     against one service may consume shared network and mitigation resources,
%     causing collateral degradation for other legitimate services.}
%     \label{fig:ddos-shared-resource-dcn}
% \end{figure}
% %---------------------------------------------------------------------

\subsection{DDoS Attack Taxonomy}
\label{subsec:ddos-taxonomy}

DDoS attacks can be grouped into several major families according to the type
of resource they attempt to exhaust. A classical taxonomy distinguishes attacks
by target resource, traffic behaviour, degree of automation, source dynamics,
and defense mechanism requirements \cite{mirkovic2004taxonomy}. This taxonomy
is important because different attack families produce different traffic
signatures and may require different mitigation actions.

\textbf{Volumetric attacks} attempt to overwhelm the available bandwidth or
packet-processing capacity of the victim network. Examples include UDP floods,
ICMP floods, and reflection or amplification attacks. These attacks are often
characterized by high packet rates, high byte rates, abnormal protocol
distributions, and sudden changes in traffic volume.

\textbf{Protocol-exhaustion attacks} exploit finite protocol state in servers,
firewalls, load balancers, or connection-tracking systems. A common example is
a SYN flood, where many half-open TCP connections are created to exhaust the
server's connection table or intermediate network devices. These attacks may
not always require the same bandwidth as volumetric attacks, but they can
cause severe degradation by exhausting stateful resources.

\textbf{Application-layer attacks} target application logic or server-side
processing resources. HTTP floods and slow-rate attacks are examples of this
category. They are more difficult to detect because they may resemble
legitimate user behavior at the network layer. Instead of relying only on raw
packet volume, detection systems may need to observe request patterns, timing
behavior, flow duration, and statistical deviations from normal traffic.

In this thesis, DDoS taxonomy is not only used for describing attacks. It also
defines the operational label space of the AI pipeline and influences how the
policy layer selects mitigation actions. For example, a UDP flood may be
mitigated by rate-limiting or filtering based on transport features, while a
slow-rate attack may require application-aware inspection or stricter
connection handling.

\subsection{DDoS Impact on Availability and SLA}
\label{subsec:ddos-sla-impact}

The operational impact of DDoS attacks can be measured through several service
quality indicators. These include throughput degradation, increased latency,
packet loss, request failure rate, connection timeout rate, and reduced service
availability. In shared-resource DCN environments, the impact should also be
evaluated from the perspective of legitimate services or tenants that are not
direct attack targets. For example, if a mitigation function consumes excessive
CPU or bandwidth while protecting one service, other services may still
experience degraded performance.

This observation motivates the SLA-aware design of \padonap{}. A DDoS defense
system should not only detect attacks accurately, but also respond in a way
that preserves legitimate service quality where resources are available. This
requires coordination between AI analytics, policy decisions, orchestration
actions, and resource allocation.

%----------------------------------------------------------------------
%----------------------------------------------------------------------
\section{CICDDoS2019 and Window-Based Traffic Representation}
\label{sec:cic-bg}
%----------------------------------------------------------------------

Machine-learning models for DDoS detection require labelled traffic data that
captures both benign traffic and diverse attack behaviours. CICDDoS2019 is a
widely used benchmark dataset for evaluating DDoS detection and classification
approaches. It was introduced to address limitations of earlier intrusion
detection datasets by providing modern DDoS attack scenarios, a DDoS taxonomy,
flow-based traffic records, and feature-based evaluation support
\cite{sharafaldin2019cicddos,cicddos2019_official}. The dataset includes
multiple DDoS attack families, including reflection/amplification attacks,
UDP-based floods, SYN-based attacks, and WebDDoS scenarios, making it suitable
for evaluating multi-class DDoS classification in the context of this thesis.

Instead of training directly on raw packets, many DDoS detection systems use
flow-level or window-level statistical representations. This design is common
in network intrusion detection because aggregated features can summarize
traffic volume, protocol behaviour, address and port distributions, connection
patterns, and timing characteristics while reducing the cost of processing
individual packets. CICDDoS2019 follows this flow-based design: raw traffic is
captured as PCAP data and processed into CSV files using CICFlowMeter, which
extracts statistical flow features such as flow duration, packet counts, byte
counts, and packet-length statistics \cite{cicddos2019_official,cicflowmeter_official}.

Window-based representation is particularly relevant to \padonap{} because the
closed-loop orchestration layer does not act on every packet individually.
Instead, it consumes telemetry summaries and inference events over recent time
intervals. A window-level representation therefore provides a bridge between
traffic observation, AI inference, policy decision making, and orchestration
actions. This design is also consistent with recent DDoS detection and
mitigation surveys, which show that statistical traffic features and
machine-learning methods remain widely used in SDN and automated defense
systems \cite{hill2024ddos_sdn_survey,apostu2025ai_ddos_survey}.

% %---------------------- FIGURE PLACEHOLDER ----------------------------
% \begin{figure}[htbp]
%     \centering
%     % \includegraphics[width=0.95\linewidth]{figures/window_feature_extraction.pdf}
%     \fbox{\parbox{0.9\linewidth}{
%     \centering
%     Placeholder for Figure: Window-based feature extraction pipeline.
%     Show packet or flow records grouped into fixed-size windows, then transformed
%     into statistical features, labels, and AI model inputs.
%     }}
%     \caption{Window-based feature extraction for DDoS analytics. Raw traffic or
%     flow records are aggregated into time windows and transformed into statistical
%     features for detection and forecasting.}
%     \label{fig:window-feature-extraction}
% \end{figure}
% %---------------------------------------------------------------------

\subsection{Window-Level Feature Groups}
\label{subsec:window-features}

The AI pipeline implemented in this thesis uses a compact set of window-level
features derived from flow records. These features are not introduced as a new
public dataset or a new feature-extraction standard. Rather, they are selected
to represent feature families commonly used in flow-based DDoS detection:
traffic rate, entropy, protocol distribution, TCP flag behaviour, packet-size
statistics, and timing behaviour. These feature families are motivated by the
observation that DDoS attacks may affect multiple measurable properties of
traffic, including volume, source and destination diversity, transport
protocol composition, connection-state indicators, and temporal regularity
\cite{mirkovic2004taxonomy,hill2024ddos_sdn_survey,cicflowmeter_official}.

\begin{table}[htbp]
\centering
\caption{Window-level feature groups used by the AI pipeline in this thesis.}
\label{tab:window-feature-groups}
\begin{tabular}{p{0.28\linewidth} p{0.62\linewidth}}
\hline
\textbf{Feature group} & \textbf{Features} \\
\hline
Traffic rate &
\texttt{pkt\_rate}, \texttt{byte\_rate}, \texttt{new\_flows\_rate} \\

Entropy features &
\texttt{src\_ip\_entropy}, \texttt{dst\_ip\_entropy},
\texttt{src\_port\_entropy}, \texttt{dst\_port\_entropy} \\

Protocol distribution &
\texttt{proto\_dist\_tcp}, \texttt{proto\_dist\_udp},
\texttt{proto\_dist\_icmp} \\

Flag and packet statistics &
\texttt{syn\_ratio}, \texttt{fin\_ratio},
\texttt{avg\_pkt\_size}, \texttt{pkt\_size\_std} \\

Timing features &
\texttt{flow\_duration\_mean}, \texttt{inter\_arrival\_mean},
\texttt{inter\_arrival\_std} \\
\hline
\end{tabular}
\end{table}

Traffic-rate features are useful for identifying sudden increases in packets,
bytes, or newly created flows. Entropy features describe whether traffic is
distributed across many sources and ports or concentrated around a smaller
number of endpoints. Protocol-distribution features capture whether traffic is
dominated by TCP, UDP, or ICMP. TCP flag statistics, such as the SYN ratio, are
useful for identifying state-exhaustion behaviour. Timing and packet-size
statistics help distinguish normal traffic from bursty, repetitive, or
slow-rate behaviour.

These feature groups follow the general principle of flow-based DDoS
detection: rather than relying on a single indicator, the model observes
multiple statistical properties of traffic over time. This is important because
different DDoS families may stress different parts of the network or service
stack. For example, volumetric attacks may be reflected in traffic-rate
features, SYN floods in TCP flag behaviour, reflection/amplification attacks
in source or destination diversity, and slow-rate attacks in timing-related
features.

\subsection{Traffic Class Label Space}
\label{subsec:traffic-label-space}

The implemented Track~A detector uses a CICDDoS2019 label grouping derived
from the attack labels available in the experiment files. The initial mapping
contains seven macro labels, but the
\texttt{WebDDoS} class produces no valid windows after the window-construction
step and is therefore removed before training. The final detector is trained
and evaluated on six classes: \texttt{BENIGN}, \texttt{Amplification},
\texttt{Syn}, \texttt{UDP-lag}, \texttt{TFTP}, and \texttt{Portmap}. This
label space follows the attack categories available in the CICDDoS2019 files
used by the experiment rather than a universal DDoS taxonomy
\cite{sharafaldin2019cicddos}.

\begin{table}[htbp]
\centering
\caption{Traffic classes used by the implemented Track~A detector.}
\label{tab:traffic-classes}
\begin{tabular}{clp{0.58\linewidth}}
\hline
\textbf{Class ID} & \textbf{Class name} & \textbf{Description} \\
\hline
0 & \texttt{BENIGN} &
Benign CICDDoS2019 traffic without an active DDoS condition. \\

1 & \texttt{Amplification} &
Reflection or amplification-style traffic. The implementation groups
\texttt{DrDoS\_*} reflection labels and the raw \texttt{UDP} label into this
macro class. \\

2 & \texttt{Syn} &
TCP SYN-based attack traffic from the CICDDoS2019 \texttt{Syn} files. \\

3 & \texttt{UDP-lag} &
UDP-lag attack traffic, retained as a minority class and oversampled during
training. \\

4 & \texttt{TFTP} &
High-volume TFTP attack traffic, kept as its own class because it dominates a
large portion of the CICDDoS2019 capture. \\

5 & \texttt{Portmap} &
Portmap attack traffic, retained as a minority class and oversampled during
training. \\
\hline
\end{tabular}
\end{table}

This label space supports immediate multi-class detection. Track~B uses a
separate binary DDoS/benign target on CICIDS2017 for multi-horizon
forecasting, so the forecasting labels are not the same as the Track~A
multi-class labels. The policy layer can still use the current predicted class
and confidence to select an appropriate response tier.

The class set is the operational label space of this thesis implementation,
not a universal DDoS taxonomy. If the system is retrained with another dataset
or deployed in another environment, the label space may need to be adjusted to
match the available attack categories and operational mitigation actions.

\subsection{Exploratory Data Analysis of CICDDoS2019}
\label{subsec:eda-cicddos2019}

Before training Track~A on CICDDoS2019, the dataset is audited to characterize
class balance, attack distribution, and feature dynamics. The full CICDDoS2019
collection contains tens of millions of flow records spread over twelve
attack sub-types plus benign traffic. To keep the operational label space
tractable, the reflection sub-types (\texttt{DrDoS\_DNS},
\texttt{DrDoS\_LDAP}, \texttt{DrDoS\_MSSQL}, \texttt{DrDoS\_NetBIOS},
\texttt{DrDoS\_NTP}, \texttt{DrDoS\_SNMP}, \texttt{DrDoS\_SSDP},
\texttt{DrDoS\_UDP}) and corresponding non-\texttt{DrDoS} reflection files
(\texttt{LDAP}, \texttt{MSSQL}, \texttt{NetBIOS}, \texttt{NTP},
\texttt{SNMP}, \texttt{SSDP}, and \texttt{UDP}) are collapsed into a single
\texttt{Amplification} macro class. The remaining explicit classes are
\texttt{Syn}, \texttt{UDP-lag}, \texttt{WebDDoS}, \texttt{TFTP}, and
\texttt{Portmap}. After feature-window construction, \texttt{WebDDoS} yields
zero valid windows and is removed, leaving the six-class detector reported in
Chapter~6.

\begin{figure}[htbp]
    \centering
    \includegraphics[width=0.95\linewidth]{thesis/figs/ai/eda_cicddos2019_class_dist.png}
    \caption{Flow-level class distribution of CICDDoS2019 after the Track~A
    macro-class grouping. Counts are shown on a log scale
    because the dataset is highly skewed: \texttt{Amplification} and
    \texttt{TFTP} dominate, while \texttt{WebDDoS} contains only a few hundred
    flows. The \texttt{BENIGN} class is intentionally rare in CICDDoS2019
    because the dataset is designed for attack-class taxonomy rather than
    benign/attack ratio realism.}
    \label{fig:eda-cicddos2019-class-dist}
\end{figure}

Three properties of CICDDoS2019 directly shape the design of Track~A. First,
the dataset is attack-dominated: more than 99\% of flows belong to an attack
class. This makes CICDDoS2019 a strong benchmark for multi-class attack
taxonomy but a poor candidate for binary forecasting that requires
benign$\rightarrow$attack transitions; the latter motivated the separate
forecasting dataset discussed in Section~\ref{subsec:eda-cicids2017}. Second,
the dataset exhibits severe per-class imbalance --- \texttt{Amplification} and
\texttt{TFTP} alone account for tens of millions of flows, while
\texttt{WebDDoS} contains a few hundred raw rows and zero valid windows after
windowing. The Track~A pipeline therefore caps
each class at approximately $25{,}000$--$50{,}000$ windows using per-class quotas
(Gaur~\&~Kumar, 2024) and applies SMOTETomek
(Batista~\textit{et~al.}, 2004; Mihoub~\textit{et~al.}, 2022) only to true
minority attack classes such as \texttt{UDP-lag} and \texttt{Portmap}.
\texttt{BENIGN} is kept unsynthesized to preserve operational distribution.
Third, each CSV file in CICDDoS2019 contains a single attack type. A naive
group-split by file would push entire classes to the test side and destroy
class balance, so Track~A uses a per-file temporal $80/20$ split with a
20-window purge gap. These three measures (macro grouping, per-class quota,
and per-file temporal split with purge) together ensure that the offline detector metrics
reported in Chapter~6 are not inflated by trivial class shortcuts.

%----------------------------------------------------------------------
\section{CICIDS2017 and Multi-Horizon Forecasting Suitability}
\label{sec:cicids2017-bg}
%----------------------------------------------------------------------

The forecasting task in Track~B requires a dataset that captures
benign$\rightarrow$attack transitions in chronological order so that a
model can learn to anticipate attacks rather than detect ongoing ones.
CICIDS2017 \cite{sharafaldin2018cicids,cicids2017_official} is adopted because it
satisfies this requirement and because it is the most-used DDoS forecasting
benchmark in the systematic review of Mittal \textit{et al.}~(2022) (29\% of
surveyed studies). CICIDS2017 is also the dataset on which the contemporary
SOTA early-warning system PARD-SSM~\cite{hiremath2026pardssm} demonstrates an
$\approx$8-minute Early Detection Margin, providing a direct point of
comparison for the 600-second strategic horizon evaluated in this thesis.

\subsection{Capture Schedule and Attack Diversity}
\label{subsec:cicids2017-schedule}

CICIDS2017 was captured over five working days at the Canadian Institute for
Cybersecurity. Monday is the benign-only reference day. The remaining four
days contain a planned mix of benign background traffic and scheduled attack
windows, including FTP/SSH brute-force, slow-rate DoS, volumetric DDoS, web
attacks, infiltration, botnet, and port scanning. The official UNB schedule
is reproduced in Table~\ref{tab:cicids2017-schedule}.

\begin{table}[htbp]
\centering
\small
\caption{Official CICIDS2017 attack schedule (Sharafaldin \textit{et al.},
2018a, Table~2). Only attack types relevant to the Track~B DDoS-class label
are highlighted in bold.}
\label{tab:cicids2017-schedule}
\begin{tabular}{@{}lllp{0.40\linewidth}@{}}
\toprule
Day & Time window & Attack & Notes \\
\midrule
Mon 03/07 & whole day      & ---                       & Benign reference \\
Tue 04/07 & 09:20 -- 10:20 & FTP-Patator               & dropped from Track~B \\
Tue 04/07 & 14:00 -- 15:00 & SSH-Patator               & dropped from Track~B \\
Wed 05/07 & 09:47 -- 10:10 & \textbf{DoS Slowloris}    & 23 min, slow-rate \\
Wed 05/07 & 10:14 -- 10:35 & \textbf{DoS Slowhttptest} & 21 min, slow-rate \\
Wed 05/07 & 10:43 -- 11:00 & \textbf{DoS Hulk}         & 17 min, volumetric \\
Wed 05/07 & 11:10 -- 11:23 & \textbf{DoS GoldenEye}    & 13 min, application flood \\
Thu 06/07 & 09:20 -- 10:42 & Web Brute Force / XSS / SQL & dropped from Track~B \\
Thu 06/07 & 14:19 -- 15:45 & Infiltration              & dropped from Track~B \\
Fri 07/07 & 10:02 -- 11:02 & Botnet ARES               & dropped from Track~B \\
Fri 07/07 & 15:56 -- 16:16 & \textbf{DDoS LOIT}        & 20 min, volumetric \\
Fri 07/07 & 13:55 -- 15:29 & Port Scan                 & dropped from Track~B \\
\bottomrule
\end{tabular}
\end{table}

\subsection{Exploratory Data Analysis of CICIDS2017}
\label{subsec:eda-cicids2017}

The Track~B data audit confirms the expected label inventory and the heavy
class imbalance.
Figure~\ref{fig:eda-cicids2017-class-dist} visualizes the flow-level label
distribution on a logarithmic axis, colour-coding each label by its treatment
in Track~B: green for benign (class~0), red for DDoS-class attacks kept as
class~1, and grey for non-DDoS attacks that are dropped at training time
rather than relabelled as benign --- this prevents contamination of the
negative class with non-DDoS attack patterns that Track~A would already
classify as a separate threat.

\begin{figure}[htbp]
    \centering
    \includegraphics[width=0.96\linewidth]{thesis/figs/ai/eda_cicids2017_class_dist.png}
    \caption{Flow-level label distribution of CICIDS2017 (counts on a log
    scale). DDoS-class attacks kept in Track~B (\texttt{DoS Hulk},
    \texttt{DDoS}, \texttt{DoS GoldenEye}, \texttt{DoS Slowloris},
    \texttt{DoS Slowhttptest}) sum to $\approx$380~K flows ($\approx$14\% of
    the kept dataset). Non-DDoS attacks shown in grey are dropped, leaving
    a clean benign/DDoS binary task.}
    \label{fig:eda-cicids2017-class-dist}
\end{figure}

After dropping non-DDoS attacks, the binary distribution becomes
$2{,}273{,}097$ benign flows (85.65\%) against $380{,}688$ DDoS-class flows
(14.35\%). This $\approx$6:1 imbalance is realistic for a deployment scenario
and is intentionally retained: Track~B is trained with Focal
Loss~\cite{lin2017focal} to handle residual imbalance, while the policy layer
applies threshold calibration before enforcement.

The single most important property of CICIDS2017 for forecasting is the
Wednesday morning attack sequence: four short, distinct DoS attacks
executed within a 90-minute window with deliberate benign gaps between them.
Figure~\ref{fig:eda-cicids2017-wed-timeline} visualizes this sequence.

\begin{figure}[htbp]
    \centering
    \includegraphics[width=0.98\linewidth]{thesis/figs/ai/eda_cicids2017_wednesday_timeline.png}
    \caption{CICIDS2017 Wednesday-morning DoS attack sequence: four short
    attacks (13--23 minutes each) interleaved by 4--10 minute benign gaps.
    This produces five distinct onset transitions and five distinct offset
    transitions within a single capture day, exactly the kind of structure
    a multi-horizon forecaster needs to be evaluated against.}
    \label{fig:eda-cicids2017-wed-timeline}
\end{figure}

The Wednesday sequence yields five benign$\rightarrow$attack onset
transitions and five attack$\rightarrow$benign offset transitions across
$\approx$2~hours of real traffic; the Friday afternoon DDoS-LOIT capture adds
one more onset/offset pair. This density of state transitions is the property
that makes CICIDS2017 superior to single-burst datasets for forecasting
research.

\subsection{Why Not CICDDoS2019 or CSE-CIC-IDS2018 for Track~B?}
\label{subsec:dataset-comparison-forecast}

A natural question is why Track~B does not reuse the Track~A dataset
(CICDDoS2019) or the larger CSE-CIC-IDS2018 capture. The empirical evidence
is captured in Figure~\ref{fig:eda-persistence-dataset-comparison}: a
persistence baseline ($\hat{y}_{t+h}=y_t$) achieves AUC~$\approx 0.95$ on
CICDDoS2019 and AUC~$\approx 0.97$ on CSE-CIC-IDS2018, while only reaching
AUC~$\approx 0.86$ on CICIDS2017. The first two datasets are essentially
single-state-per-day: CICDDoS2019 attack days are $>99\%$ attack flows and
CSE-CIC-IDS2018 contains one multi-hour attack burst per day. A trivial
copy-the-current-second forecaster therefore wins almost automatically,
leaving no headroom for a deep model to demonstrate learning. CICIDS2017
breaks this trivial baseline because the Wednesday DoS sequence forces a
forecaster to handle five onset/offset transitions per day.

\begin{figure}[htbp]
    \centering
    \includegraphics[width=0.92\linewidth]{thesis/figs/ai/eda_persistence_dataset_comparison.png}
    \caption{Persistence-baseline AUC as a dataset-difficulty proxy for
    multi-horizon forecasting. CICDDoS2019 and CSE-CIC-IDS2018 produce
    near-trivial persistence performance because attack episodes dominate
    the capture, while CICIDS2017's mixed Wednesday DoS sequence drops
    persistence to $\approx 0.86$, opening a measurable headroom for learned
    models. This empirical comparison motivated the dataset switch to
    CICIDS2017 documented in Chapter~6.}
    \label{fig:eda-persistence-dataset-comparison}
\end{figure}

In summary, Track~B adopts CICIDS2017 because (i) it is the most-cited
forecasting benchmark in the field, (ii) it is the dataset on which the SOTA
PARD-SSM~\cite{hiremath2026pardssm} reports its eight-minute Early Detection
Margin, and (iii) its Wednesday DoS sequence provides the genuine onset
transitions required to evaluate proactive defense.

%----------------------------------------------------------------------
\section{SDN, NFV, CNFs, and Service Function Chaining}
\label{sec:nfv-bg}
%----------------------------------------------------------------------

Dynamic DDoS mitigation requires the ability to observe traffic, decide on a
response, and modify the network or service path accordingly. Software-Defined
Networking, Network Function Virtualization, Cloud-Native Network Functions,
and Service Function Chaining provide the architectural foundations for this
type of programmable defense.

\subsection{Software-Defined Networking}
\label{subsec:sdn-bg}
\begin{figure}[htbp]
    \centering
      \includegraphics[width=\textwidth]{thesis/figs/SDN.png}%
    \caption{Overview of SDN architecture.}
    \label{fig:arch}
\end{figure}

Software-Defined Networking (SDN) separates the control plane from the data
plane, allowing forwarding behavior to be controlled through a logically
centralized controller \cite{sdn_survey_kreutz2015}. In traditional networks,
routing and forwarding decisions are distributed across individual devices,
which makes global policy enforcement difficult. SDN introduces a programmable
control layer that can install forwarding rules, redirect traffic, and enforce
network-wide policies.

A typical SDN architecture can be viewed as three logical layers. The
application layer contains network applications such as routing, monitoring,
traffic engineering, and security policy enforcement. The control layer hosts
the SDN controller, which maintains a global view of the network and translates
high-level policy intent into forwarding decisions. The infrastructure layer
contains data-plane devices, such as switches and routers, that forward packets
according to the rules installed by the controller.

This separation is important for DDoS defense because mitigation often requires
fast changes to the forwarding path rather than only offline attack detection.
Once an attack is detected or forecasted, the control logic can install or
update forwarding rules that redirect suspicious traffic through a mitigation
function, such as a rate limiter or scrubber, while allowing legitimate traffic
to continue toward the protected service. In the scope of this thesis, SDN is
therefore treated as the programmable traffic-steering foundation for
service-function insertion rather than as the main research contribution.

In the \padonap{} testbed, this SDN principle is represented by Open vSwitch
rule installation and service-function chaining. The policy and orchestration
logic selects a response tier, and the enforcement layer updates the forwarding
behavior so that selected flows traverse the appropriate mitigation function.

\subsection{Network Function Virtualization and CNFs}
\label{subsec:nfv-cnf-bg}

Network Function Virtualization (NFV) decouples network functions from
specialized hardware and allows them to run as software instances on commodity
infrastructure \cite{etsi_nfv_architecture}. Examples of virtualized network
functions include firewalls, intrusion detection systems, load balancers, rate
limiters, and DDoS scrubbers. NFV improves flexibility because operators can
instantiate, scale, or remove functions without deploying dedicated hardware
appliances.

Cloud-Native Network Functions (CNFs) extend this idea by packaging network
functions as containers managed by platforms such as Kubernetes
\cite{kubernetes_docs}. Compared with VM-based VNFs, CNFs are often easier to
package, deploy, update, and observe using cloud-native deployment technologies
such as Kubernetes. However, CNFs are not instantly enforceable. A CNF must be
scheduled, started, marked ready, and connected to the traffic path before it
can affect live traffic. This readiness delay is a central concern in
proactive DDoS mitigation.

\subsection{Service Function Chaining}
\label{subsec:sfc-bg}

Service Function Chaining (SFC) describes the ordered steering of traffic
through a sequence of service functions \cite{sfc_rfc7665}. A simple service
chain for DDoS defense may redirect traffic through a firewall, rate limiter,
or scrubber before forwarding it to the protected application.

% %---------------------- FIGURE PLACEHOLDER ----------------------------
% \begin{figure}[htbp]
%     \centering
%     % \includegraphics[width=0.95\linewidth]{figures/sfc_ddos_mitigation.pdf}
%     \fbox{\parbox{0.9\linewidth}{
%     \centering
%     Placeholder for Figure: Service function chaining for DDoS mitigation.
%     Show client or attack traffic passing through firewall, rate limiter,
%     scrubber, and then protected service.
%     }}
%     \caption{Example service function chain for DDoS mitigation. Traffic may be
%     steered through filtering, rate-limiting, or scrubbing functions before
%     reaching the protected service.}
%     \label{fig:sfc-ddos-mitigation}
% \end{figure}
% %---------------------------------------------------------------------

In \padonap{}, SFC denotes the traffic-steering behavior that connects the
protected service to the selected mitigation function. When OVS rules are
actually installed in the Mininet testbed, the evaluation describes this as
implemented SFC-style enforcement. When a lifecycle step is represented only
by the orchestration stub, it is reported as emulated behavior.

\subsection{Mitigation Delay in NFV/CNF Environments}
\label{subsec:mitigation-delay-bg}

Although NFV and CNF technologies make mitigation functions easier to deploy
and scale, they do not eliminate the delay between attack observation and
effective enforcement. A purely reactive defense loop must wait until
suspicious traffic is observed, classified, evaluated by policy, sent to the
orchestrator, deployed as a mitigation function, and connected to the service
path. This interval is referred to in this thesis as the
\attack-to-enforcement latency or mitigation exposure window.

A simplified latency decomposition is:

\[
T_{\text{enforce}} =
T_{\text{collect}} +
T_{\text{detect}} +
T_{\text{policy}} +
T_{\text{orchestrate}} +
T_{\text{deploy}} +
T_{\text{ready}} +
T_{\text{steer}}.
\]

Here, \(T_{\text{collect}}\) denotes telemetry collection latency,
\(T_{\text{detect}}\) denotes AI inference latency, \(T_{\text{policy}}\)
denotes policy decision latency, \(T_{\text{orchestrate}}\) denotes the time
needed to submit and process an orchestration request, \(T_{\text{deploy}}\)
denotes CNF deployment time, \(T_{\text{ready}}\) denotes readiness-check time,
and \(T_{\text{steer}}\) denotes the time required to apply traffic steering or
filtering.

% %---------------------- FIGURE PLACEHOLDER ----------------------------
% \begin{figure}[htbp]
%     \centering
%     % \includegraphics[width=0.95\linewidth]{figures/reactive_mitigation_timeline.pdf}
%     \fbox{\parbox{0.9\linewidth}{
%     \centering
%     Placeholder for Figure: Reactive mitigation timeline.
%     Show attack start, telemetry collection, detection, policy decision,
%     orchestration, CNF deployment, readiness, traffic steering, and mitigation
%     active. Highlight the exposure window.
%     }}
%     \caption{Reactive mitigation timeline in NFV/CNF environments. The exposure
%     window includes telemetry collection, detection, policy decision,
%     orchestration, CNF readiness, and enforcement delays.}
%     \label{fig:reactive-mitigation-timeline}
% \end{figure}
% %---------------------------------------------------------------------

During this exposure window, malicious traffic may continue to consume shared
resources and degrade service quality. The purpose of multi-horizon forecasting
in \padonap{} is therefore not only to predict future attack states, but also
to create useful lead time for preparing mitigation before the attack reaches
its most damaging stage.

%----------------------------------------------------------------------
\section{ONAP}
\label{sec:onap-bg}
%----------------------------------------------------------------------

\begin{figure}[htbp]
    \centering
      \includegraphics[width=\textwidth]{thesis/figs/ONAP.png}%
    \caption{Overview of the ONAP High-Level 
    architecture.}
    \label{fig:arch}
\end{figure}
The Open Network Automation Platform (ONAP) is an open-source platform for
network automation and service lifecycle management \cite{onap_docs}. ONAP is
designed to support model-driven service design, orchestration, monitoring,
policy evaluation, and closed-loop control. Although a full ONAP deployment is
large and complex, its architecture provides useful concepts for structuring
automated mitigation systems.

\subsection{ONAP Components Relevant to \padonap{}}
\label{subsec:onap-components-bg}


Several ONAP concepts are especially relevant to \padonap{}.

\textbf{DCAE} collects telemetry, analytics outputs, and event data from the
managed environment. In a closed-loop system, this
collection function provides the event stream that triggers analysis and
policy decisions.

\textbf{The Policy Framework} evaluates policies and determines the appropriate
operational response. In the context of DDoS
mitigation, policy logic may decide whether to remain in normal monitoring
mode, raise an alert, pre-position a mitigation CNF, activate rate limiting, or
isolate a suspicious flow or tenant.

\textbf{The Service Orchestrator} supports service and network-function
lifecycle operations. For \padonap{}, the SO-like role is
to translate policy decisions into deployment or update actions in the
Kubernetes environment.

Together, these components define a closed-loop automation path in which
network events are observed, analyzed, converted into policy decisions, and
translated into operational actions.

\begin{table}[htbp]
\centering
\small
\caption{ONAP concepts adopted by \padonap{} and their implementation boundary.}
\label{tab:onap-concepts-bg}
\begin{tabular}{@{}p{0.20\linewidth}p{0.36\linewidth}p{0.34\linewidth}@{}}
\toprule
ONAP concept & Role in closed-loop automation & Role in \padonap{} \\
\midrule
DCAE & Event collection and analytics ingestion &
DCAE-like collector receives AI output events and forwards them to policy
logic. \\
Policy Framework & Policy evaluation and action selection &
Policy-like engine maps detection, forecast, SHAP metadata, and resource state
to response tiers. \\
SO & Service lifecycle orchestration &
SO-like adapter translates tier decisions into CNF deployment, scaling, or
update requests. \\
A\&AI & Inventory and topology state &
Represented as lightweight topology and resource metadata in the testbed. \\
\bottomrule
\end{tabular}
\end{table}

\subsection{ONAP-Inspired in \padonap{}}
\label{subsec:onap-inspired-bg}

\padonap{} follows the functional structure of ONAP automation,
but it is treated as an ONAP-inspired implementation rather than a full
ONAP deployment except where the selected ONAP OOM experiment is explicitly
reported.

The control path adopted in this thesis is:

\begin{figure}[htbp]
\centering
\begin{tikzpicture}[
    node distance=1.25cm,
    >=Latex,
    block/.style={
        rectangle,
        rounded corners=3pt,
        draw=black!55,
        very thick,
        align=center,
        minimum width=2.65cm,
        minimum height=1.05cm,
        fill=blue!6,
        blur shadow
    },
    action/.style={
        rectangle,
        rounded corners=3pt,
        draw=black!55,
        very thick,
        align=center,
        minimum width=2.85cm,
        minimum height=1.05cm,
        fill=green!8,
        blur shadow
    },
    arrow/.style={
        ->,
        very thick,
        draw=black!70
    }
]

\node[block] (telemetry) {Telemetry\\Sources};
\node[block, right=of telemetry] (ai) {AI\\Inference};
\node[block, right=of ai] (dcae) {DCAE-like\\Collector};
\node[action, below=1.25cm of dcae] (policy) {Policy\\Engine};
\node[action, left=of policy] (so) {SO-like\\Adapter};
\node[action, left=of so] (k8s) {Kubernetes/CNF\\Enforcement};

\draw[arrow] (telemetry) -- (ai);
\draw[arrow] (ai) -- (dcae);
\draw[arrow] (dcae) -- (policy);
\draw[arrow] (policy) -- (so);
\draw[arrow] (so) -- (k8s);

\draw[arrow, dashed, bend left=35] (k8s.north) to node[above, align=center, font=\small]
{updated mitigation state} (telemetry.south);

\node[
    draw=black!35,
    rounded corners=4pt,
    fit=(telemetry)(ai)(dcae)(policy)(so)(k8s),
    inner sep=0.35cm,
    label={[font=\bfseries]above:ONAP-inspired \padonap{} Closed Loop}
] {};

\end{tikzpicture}
\caption{ONAP-inspired closed-loop control path in \padonap{}. Telemetry is
processed by the AI inference engine, collected by a DCAE-like component,
translated into policy decisions, and executed through an SO-like adapter and
Kubernetes-based CNF enforcement.}
\label{fig:padom-closed-loop-tikz}
\end{figure}

This design allows the thesis to study the operational value of AI-assisted
closed-loop mitigation without requiring the full complexity of production
ONAP. The DCAE-like collector receives AI inference events, the policy engine
maps these events to response tiers, the SO-like adapter performs orchestration
actions, and Kubernetes enforces the resulting CNF deployment or update.

%----------------------------------------------------------------------
\section{Streaming Telemetry with gNMI, Kafka, and Flink}
\label{sec:streaming-telemetry-bg}
%----------------------------------------------------------------------

Closed-loop DDoS mitigation has a hard real-time component: telemetry must be
collected, transported, and aggregated quickly enough that the AI inference,
policy decision, and orchestration enforcement still complete inside the
exposure window of the attack. Classical pull-based network-management
mechanisms such as SNMP \texttt{GET} polling are ill-suited to this role
because they introduce two undesirable couplings. First, the polling
frequency is bounded below by the manager's request RTT, so increasing the
sampling rate proportionally inflates the management-plane load on every
network device. Second, every poll is a separate request--response
transaction, so the device state observed by the manager is always at least
one RTT behind the device's actual state, and there is no built-in way to
buffer or replay missed samples. Streaming telemetry inverts both couplings:
the device pushes structured updates on a subscription, and a downstream
message bus absorbs short bursts so the analytics path can run at its own
pace. In \padonap{}, this streaming role is implemented by a three-stage
pipeline summarized in
Figure~\ref{fig:streaming-telemetry-overview} and detailed in the next three
subsections: (i)~\textbf{gNMI/gRPC} collection at the data-plane edge,
(ii)~\textbf{Apache Kafka} message transport, and (iii)~\textbf{Apache Flink}
stateful stream processing.

\begin{figure}[H]
\centering
\resizebox{\textwidth}{!}{\begin{tikzpicture}[
  font=\sffamily\scriptsize,
  >=Latex,
  box/.style={draw, rounded corners=2pt, align=center, inner sep=4pt,
              minimum height=10mm, minimum width=22mm},
  src/.style={box, fill=blue!8,  draw=blue!55!black},
  bus/.style={box, fill=orange!10, draw=orange!70!black},
  proc/.style={box, fill=green!10, draw=green!55!black},
  sink/.style={box, fill=purple!10, draw=purple!55!black},
  flow/.style={->, thick}
]
\node[src]  (sw)     at (0,0)    {OVS / vendor NOS\\(\texttt{ifctr}, flow ctrs)};
\node[src]  (gnmi)   at (3.6,0)  {gNMI server /\\compat.\ adapter};
\node[bus]  (kafka)  at (7.4,0)  {Kafka topic\\\texttt{pad.telemetry.raw}};
\node[proc] (flink)  at (11.0,0) {Flink job\\5\,s sliding window};
\node[sink] (ai)     at (14.6,0) {AI feature\\vector $\mathbf{x}_t$};
\draw[flow] (sw)    -- node[above, font=\sffamily\tiny]{counters / NetFlow v5}  (gnmi);
\draw[flow] (gnmi)  -- node[above, font=\sffamily\tiny]{\texttt{Subscribe\,(STREAM)}} (kafka);
\draw[flow] (kafka) -- node[above, font=\sffamily\tiny]{poll consumer}          (flink);
\draw[flow] (flink) -- node[above, font=\sffamily\tiny]{17 features /\,1 s}    (ai);

\draw[dashed, gray] (-1.0,-1.2) -- (14.6+1.6,-1.2);
\node[font=\sffamily\tiny, gray] at (1.8,-1.45) {Data plane};
\node[font=\sffamily\tiny, gray] at (5.5,-1.45) {Collection};
\node[font=\sffamily\tiny, gray] at (9.2,-1.45) {Transport};
\node[font=\sffamily\tiny, gray] at (12.8,-1.45) {Stream processing};
\end{tikzpicture}}
\caption{End-to-end streaming-telemetry pipeline used by \padonap{}: switch
counters or NetFlow v5 records are exported by a gNMI server (or a
gNMI-compatible adapter) to a Kafka topic, then aggregated into per-window
feature vectors by Apache Flink. Only the Kafka topic crosses the
local$\leftrightarrow$remote boundary in the evaluated deployment.}
\label{fig:streaming-telemetry-overview}
\end{figure}

\subsection{gNMI/gRPC Model-Driven Telemetry}
\label{subsec:gnmi-bg}

The gRPC Network Management Interface (gNMI) is a model-driven RPC interface
for retrieving and streaming network state through structured paths and typed
updates \cite{gnmi_spec}. A gNMI server exposes a tree of operational data
identified by paths (e.g.\
\path{/interfaces/interface[name=eth0]/state/counters/in-pkts}); the schema
of every path is described by a YANG module
\cite{yang_rfc7950}, so each value carries an explicit type, unit, and
metadata. Communication runs over gRPC \cite{grpc_docs}, which gives gNMI two
properties relevant to a streaming pipeline: a binary protobuf encoding that
is compact enough for high-rate counters, and HTTP/2 multiplexing so that
multiple long-lived subscriptions share a single TCP connection.

gNMI provides four RPC methods --- \texttt{Capabilities}, \texttt{Get},
\texttt{Set}, and \texttt{Subscribe} --- of which \texttt{Subscribe} is the
one relevant to closed-loop telemetry. A \texttt{SubscribeRequest} carries a
list of paths and one of three subscription modes:

\begin{itemize}
    \item \textbf{ONCE} --- the server sends one update per path and closes
          the stream; useful for snapshot debugging but not for closed-loop
          telemetry.
    \item \textbf{POLL} --- the server keeps the stream open but only emits
          updates when the client sends an explicit \texttt{Poll} message;
          retains the pull-based behaviour of SNMP while reusing the gRPC
          connection.
    \item \textbf{STREAM} --- the server proactively pushes updates either
          on change (\texttt{ON\_CHANGE}), at a fixed interval
          (\texttt{SAMPLE} with a sample period $\Delta t$), or only when
          the value moves outside a configured band
          (\texttt{TARGET\_DEFINED}). This is the mode used by \padonap{}.
\end{itemize}

For per-second AI inference, \padonap{} subscribes in \texttt{STREAM} with
\texttt{SAMPLE} and a 500\,ms sample interval. Two further options shape the
pipeline behaviour. The \texttt{suppress\_redundant} flag asks the server to
omit updates whose value has not changed since the last sample, eliminating
quiescent traffic from the bus. The \texttt{heartbeat\_interval} flag asks
the server to still emit a periodic update even when the value is suppressed,
so the downstream pipeline can distinguish ``no change'' from ``connection
silently lost''. Together, these flags give the AI pipeline a bounded
worst-case freshness while keeping bandwidth usage proportional to actual
state change.

Compared to packet-level capture, gNMI counter streams are much cheaper to
transport (one fixed-size protobuf per sample versus full PCAP frames) and
already carry the device's authoritative view of port and flow statistics.
Compared to SNMP polling, the same fields are delivered with sub-second
latency, no per-sample request RTT, and a structured schema that the
downstream Flink job can deserialize without ad-hoc parsing
\cite{su2024ddos_sdn_survey}. These properties together make gNMI/gRPC a
natural front-end for a Kafka-fed AI pipeline.

\subsection{Apache Kafka as the Telemetry Bus}
\label{subsec:kafka-bg}

The next stage of the pipeline is Apache Kafka, a distributed
publish--subscribe log designed for high-throughput telemetry-style workloads
\cite{kafka_kreps2011}. Kafka organizes data as topics, where each
topic is split into one or more ordered, append-only partitions that
are replicated across a configurable set of brokers. Producers write records
into a partition by appending to the log; consumers read records by
maintaining an offset pointer into the log and advancing it as they
process events. Three properties of this model matter for streaming
telemetry:

\begin{itemize}
    \item \textbf{Producer--consumer decoupling.} The gNMI collector
          publishes to \texttt{pad.telemetry.raw} without knowing how many
          consumers exist or how fast they read. Adding an audit consumer or
          a second AI pipeline only requires a new consumer group; it does
          not slow down the producer. In \padonap{}, this is what lets the
          local Mininet collector forward records to a remote Kubernetes
          cluster without any control-loop interaction beyond the message
          itself (Figure~\ref{fig:testbed-remote-pipeline}).

    \item \textbf{Ordered partitions and parallelism.} Within a partition,
          Kafka guarantees first-in/first-out order. Because each Flink
          subtask is assigned an exclusive subset of partitions, the
          downstream window aggregation operates on per-key chronological
          streams without an explicit sort. The number of partitions is
          chosen to match the Flink job parallelism so that no partition
          becomes a bottleneck.

    \item \textbf{Retention and replay.} Each topic retains records for a
          configurable window (a few hours of telemetry in the evaluated setting). A
          consumer that falls behind --- for example, because the AI
          container restarts --- can rewind its offset and replay the buffer
          instead of dropping events. This converts a transient processing
          stall into bounded latency rather than data loss.
\end{itemize}

Kafka's delivery semantics are configurable through producer acknowledgement
and consumer commit modes. The default \texttt{acks=1} setting commits a
record after the partition leader has written it but before replication
completes, giving low producer latency at the cost of at-most-once delivery.
\padonap{} uses \texttt{acks=all} (record committed only after all in-sync
replicas acknowledge) combined with Kafka's transactional producer support
to achieve effectively-once semantics on the producer side, while the Flink
consumer uses checkpointed offsets so that consumer-side restarts do not
double-count windows.

\subsection{Apache Flink for Stateful Stream Processing}
\label{subsec:flink-bg}

The final stage of the pipeline is Apache Flink, which performs the
window-level feature aggregation that the AI pipeline consumes. Flink models
a streaming computation as a directed acyclic graph of operators connected by
data-flow edges and provides first-class support for stateful event-time
processing \cite{flink_carbone2015}. Three of its abstractions are central
to \padonap{}: time semantics, windows, and managed state.

\paragraph{Event time vs.\ processing time.}
Flink distinguishes between an event's event time (the moment the
underlying observation was generated, carried in the record itself) and its
processing time (the wall-clock moment Flink received it). For DDoS
telemetry these can diverge: the gNMI collector and the AI sink may be in
different datacenters, and the network in between can briefly slow down. To
keep window boundaries aligned with the device's clock rather than the
broker's, Flink jobs in \padonap{} use event-time semantics keyed on the
gNMI sample timestamp. Out-of-order arrivals are handled through
watermarks: Flink only finalizes a window once it has observed a
watermark $w_t$ such that $w_t \geq$ the window upper bound, with a small
configured slack to absorb expected jitter.

\paragraph{Window types.}
Flink supports three window flavours that are useful for traffic analytics:

\begin{itemize}
    \item \textbf{Tumbling windows} of fixed length and zero overlap.
          Inexpensive but coarse-grained: a 5\,s tumbling window emits one
          feature vector every 5\,s, so the AI loop wakes only five times
          per attack-onset second.
    \item \textbf{Sliding windows} of length $L$ and slide $S<L$, producing
          a new feature vector every $S$ seconds while still aggregating
          $L$ seconds of history. \padonap{} uses $L=5$\,s, $S=1$\,s,
          so the AI engine sees a 5-second moving average refreshed every
          second.
    \item \textbf{Session windows} that close after a gap of inactivity.
          Used in \padonap{} only for ad hoc flow-cluster analysis, not for
          the main inference loop.
\end{itemize}

A sliding $L{=}5$\,s, $S{=}1$\,s window matches the latency budget of the
closed loop: 1\,s slide is short enough to feed the policy engine in
near-real-time, while 5\,s aggregation is long enough to average out
sub-second noise from the gNMI sample interval.

\paragraph{Managed state and exactly-once.}
Window aggregation is stateful by nature --- every running window holds an
accumulator until it fires. Flink stores this state in a managed RocksDB or
heap backend and snapshots it through periodic checkpoints aligned
across all operators. If a Flink TaskManager crashes, the job recovers from
the latest checkpoint and replays Kafka records from the corresponding
offset, so window outputs are not duplicated and not lost. Together with
Kafka's transactional producer, this delivers an effectively
exactly-once guarantee on the AI feature stream.

\paragraph{Feature extractor pipeline.}
The Flink job implemented in this thesis is a five-operator chain:
\texttt{KafkaSource} $\rightarrow$ \texttt{KeyBy(switch\_id, port)}
$\rightarrow$ \texttt{SlidingEventTimeWindow(5\,s, 1\,s)}
$\rightarrow$ \texttt{ProcessFunction (compute 17-d feature vector)}
$\rightarrow$ \texttt{KafkaSink (\texttt{pad.features.v1})}.
\texttt{KeyBy} partitions the stream so that per-port statistics are
computed in parallel without cross-talk; the \texttt{ProcessFunction}
implements the closed-form rate, entropy, protocol, flag, and timing
features listed in Table~\ref{tab:features}. The output topic is consumed
by the XGBoost detector and the CNN+BiLSTM+Attention forecaster running
inside the remote Kubernetes cluster.

\subsection{Testbed Adaptation: gNMI-Compatible OVS Adapter and softflowd}
\label{subsec:telemetry-adaptation-bg}

The thesis testbed runs on Open vSwitch, which historically does not expose
a production-grade native gNMI server, and the production deployment
described in Section~\ref{subsec:cicids2017-schedule} additionally needs to
capture per-host NetFlow records that are not visible to the switch counter
view. Two adapter mechanisms are therefore used:

\begin{itemize}
    \item A \textbf{gNMI-compatible adapter} samples OVS port and flow
          counters every 500\,ms through the OpenFlow / OVSDB control
          interfaces, normalizes the values into a YANG-like protobuf
          record using the same field layout an upstream gNMI server would
          produce, and publishes those records to Kafka. This adapter is
          used for synthetic scenario replay where switch counters are the
          authoritative traffic source.
    \item A \textbf{softflowd-based NetFlow~v5 exporter} runs on every host
          (per Figure~\ref{fig:testbed-remote-pipeline}). It samples the
          per-flow tuples that softflowd reconstructs from packet traces,
          encodes them in NetFlow~v5, and forwards them to a Kafka
          producer in the same topic schema as the gNMI adapter. This
          adapter is used in the live Mininet run because it captures
          end-host visible traffic before any OVS-side rate limiting,
          giving a faithful view of the offered load.
\end{itemize}

From the downstream pipeline's perspective the two adapters are
interchangeable: Flink consumes the same record schema, computes the same
17-dimensional feature vector, and feeds the same AI track. The boundary is
explicit in the thesis text because the term ``gNMI-compatible'' refers to
schema and access-pattern equivalence, not to a native gNMI implementation in
OVS. This honest framing preserves the architectural value of a gNMI / Kafka
/ Flink streaming-telemetry stack while reflecting the actual capabilities
of the testbed.

%----------------------------------------------------------------------
\section{AI-Based Detection and Multi-Horizon Forecasting}
\label{sec:ai-bg}
%----------------------------------------------------------------------

AI analytics in \padonap{} has two complementary roles. The first role is
immediate detection: identifying the current traffic class and estimating
whether an attack is already occurring. The second role is forecasting:
estimating a future binary DDoS risk indicator over multiple horizons so that the
orchestration loop can prepare mitigation before the attack reaches its most
damaging stage.

\begin{figure}[htbp]
    \centering
    \includegraphics[width=0.95\linewidth]{thesis/figs/ai/dual_track_ai_engine.png}
    \caption{Dual-track AI engine in \padonap{}. Track~A (XGBoost) produces an
    immediate multi-class label and confidence; Track~B (CNN+BiLSTM+Attention)
    produces three forecast probabilities $\hat p_{t+60}, \hat p_{t+300},
    \hat p_{t+600}$. Both outputs are consumed by the policy engine for
    response-tier selection.}
    \label{fig:dual-track-ai-engine}
\end{figure}

\subsection{XGBoost-Based Immediate DDoS Classification}
\label{sec:xgb-bg}

Track A of the AI pipeline performs immediate DDoS classification. After a
controlled comparison of six classifier families (Section~\ref{subsec:classifier-comparison}),
\textbf{XGBoost} is retained as the detector used for
policy-level explanations. XGBoost is a scalable gradient-boosted tree method
that builds an ensemble of decision trees to optimize a regularized objective
function \cite{xgboost_chen2016} and is well suited to tabular network-flow
features because it can model non-linear feature interactions and support
SHAP-based explanation with a mature tree explainer.

In the AI pipeline implemented for this thesis, Track A is defined as:

\[
\text{Window-level traffic features}
\rightarrow
\text{XGBoost classifier}
\rightarrow
\text{6-class prediction and confidence}.
\]

The classifier outputs a predicted class and a confidence score. These outputs
can be consumed by the policy engine to support reactive mitigation or
escalation. For example, if the classifier predicts \texttt{Syn} with
high confidence, the policy engine may select a mitigation tier that enables
SYN-specific filtering or rate limiting.

However, AI performance is reported only where it is supported by experiment
artifacts, saved model outputs, or reproducible evaluation logs.
Unsupported values are excluded from the quantitative evaluation.

\subsection{CNN+BiLSTM+Attention Multi-Horizon Forecasting}
\label{sec:forecast-bg}

Track B of the AI pipeline addresses proactive DDoS defense through
multi-horizon forecasting. Unlike immediate classification, which estimates the
current traffic class, forecasting estimates a near-future attack indicator
given a sequence of recent traffic windows. This is useful because
orchestration and CNF readiness introduce delay before mitigation becomes
enforceable.

To select an architecture for Track~B, five sequence-modelling
families on the same lookback/horizon setup: Stacked LSTM, Gated Recurrent
Unit (GRU), Temporal Convolutional Network (TCN), Transformer encoder paired
with an LSTM head, and a CNN + Bidirectional LSTM with attention pooling
(``CNN+BiLSTM+Attention'') are benchmarked. The CNN+BiLSTM+Attention design follows recent
DDoS-forecasting literature, where a convolutional front-end extracts local
temporal patterns, a bidirectional LSTM models longer-range context, and an
attention layer learns dynamic temporal weights over the lookback window
\cite{alghazzawi2021cnnbilstm, alanazi2022ensembleddos, hamcid2025nature,
cnnbilstmattn2025}. Two sanity baselines are added to the comparison: a
persistence forecaster $\hat{y}_{t+h} = y_t$ \cite{hyndman2018forecasting} and
a tabular XGBoost trained on the flattened lookback, which removes the
sequence inductive bias while keeping the same input
\cite{zerveas2021tstransformer}.

The forecaster predicts a binary DDoS-class indicator at three operationally
motivated horizons:

\[
h \in \{60~\text{s},\ 300~\text{s},\ 600~\text{s}\}.
\]

The 60~s horizon (tactical) matches an immediate firewall or
rate-limiter response, the 300~s (5-min, operational) horizon matches
automated mitigation activation, and the 600~s (10-min, strategic)
horizon is comparable to the eight-minute Early Detection Margin reported by
PARD-SSM~\cite{hiremath2026pardssm}, the closest contemporary DDoS
early-warning baseline that uses the same dataset.

\subsection{Forecasting Lead Time and False Positive Risk}
\label{subsec:forecast-leadtime-bg}

Forecasting is valuable because it can create lead time for proactive
mitigation. For example, if the model assigns high future risk to a UDP-flood
episode, the policy engine can pre-position or pre-scale a mitigation CNF
before the traffic peak arrives. This reduces dependence on a purely reactive
threshold-crossing mechanism.

However, forecasting also introduces false positive risk. If the model
incorrectly predicts an attack during benign traffic fluctuations, the system
may waste compute resources, disturb legitimate traffic, or trigger unnecessary
mitigation actions. Therefore, \padonap{} does not treat every forecast as a
direct enforcement command. Instead, the policy layer combines forecast
confidence, detection confidence, horizon length, recent traffic behaviour, and
resource feasibility before selecting a response tier.

%----------------------------------------------------------------------
\section{Explainable Policy Translation and Response Tiers}
\label{sec:policy-bg}
%----------------------------------------------------------------------

AI models do not directly enforce mitigation actions. Their outputs must be
translated into operational decisions that can be understood and executed by
the orchestration layer. This translation step is essential because a practical
defense system must decide not only whether traffic is malicious, but also
what action should be taken, when it should be taken, and why that action is
justified.

\subsection{From AI Outputs to Policy Decisions}
\label{subsec:ai-to-policy-bg}

In \padonap{}, the policy layer consumes multiple types of input metadata:

\begin{itemize}
    \item Current predicted traffic class;
    \item Detection confidence;
    \item Forecasted future DDoS risk;
    \item Forecast horizon;
    \item Forecast confidence or risk score;
    \item Selected traffic indicators;
    \item SHAP-based feature explanations, if available;
    \item Resource state and allocation feasibility.
\end{itemize}

The policy layer maps these inputs to a response tier. This design separates
AI inference from operational enforcement. The AI model provides evidence, but
the policy layer decides whether that evidence is strong enough to justify
alerting, preparation, mitigation, or isolation.

\subsection{Response Tier Model}
\label{subsec:response-tier-bg}

A tiered response model is useful because not all warnings should trigger
immediate enforcement. Low-confidence forecasts may justify increased
monitoring, while high-confidence detection and forecast agreement may justify
active mitigation.

\begin{table}[htbp]
\centering
\caption{Example response tiers for proactive DDoS mitigation in \padonap{}.}
\label{tab:response-tiers-bg}
\begin{tabular}{p{0.14\linewidth} p{0.24\linewidth} p{0.52\linewidth}}
\hline
\textbf{Tier} & \textbf{Name} & \textbf{Purpose} \\
\hline
0 & Normal &
No mitigation action; continue telemetry collection and AI monitoring. \\

1 & Alert &
Record warning, increase observation sensitivity, and prepare policy context. \\

2 & Pre-position &
Prepare or pre-scale lightweight mitigation CNFs without applying disruptive
traffic enforcement. \\

3 & Mitigation &
Activate rate limiting, filtering, scrubbing, or traffic steering actions. \\

4 & Isolation &
Apply stronger containment for severe, persistent, or high-confidence attacks. \\
\hline
\end{tabular}
\end{table}

This tiered approach supports proactive defense while reducing the risk of
unnecessary mitigation. For example, a forecast-only warning may trigger
Tier~2 pre-positioning, while agreement between immediate detection and
short-horizon forecasting may trigger Tier~3 mitigation.

% %---------------------- FIGURE PLACEHOLDER ----------------------------
% \begin{figure}[htbp]
%     \centering
%     % \includegraphics[width=0.95\linewidth]{figures/ai_policy_translation.pdf}
%     \fbox{\parbox{0.9\linewidth}{
%     \centering
%     Placeholder for Figure: AI-to-policy translation.
%     Show detection output, forecast output, SHAP metadata, and resource state
%     entering a policy engine, which selects Tier 0--4 and sends action to the
%     orchestrator.
%     }}
%     \caption{Translation of AI analytics outputs into auditable policy decisions
%     and graduated mitigation tiers.}
%     \label{fig:ai-policy-translation}
% \end{figure}
% %---------------------------------------------------------------------

\subsection{SHAP Explainability and Audit Metadata}
\label{sec:shap-bg}

Explainability is important in security automation because operators need to
understand why an AI model recommends a response. SHAP provides feature
attribution values based on a game-theoretic formulation of feature
contribution \cite{shap_lundberg2017}. For a DDoS classifier, SHAP can show
whether a prediction was influenced by packet-rate features, entropy features,
SYN behaviour, protocol distribution, or timing statistics.

In \padonap{}, SHAP explanations can be attached to AI inference outputs and
passed to the policy engine as explanation metadata. This supports
auditability because a mitigation decision can be associated with the features
that most strongly influenced the model. Example audit metadata may include:

\begin{itemize}
    \item Timestamp of the inference event;
    \item Predicted class and confidence;
    \item Forecasted class and horizon;
    \item Selected response tier;
    \item Top contributing features;
    \item Resource feasibility result;
    \item Orchestration action taken.
\end{itemize}

SHAP is not used as an independent mitigation trigger in this thesis. The
policy tier is selected from detection confidence, forecasted risk, hysteresis,
and resource feasibility. SHAP instead provides an explanation layer: for
example, a \texttt{UDP\_Flood} decision can be accompanied by high
contributions from packet rate, UDP protocol fraction, and low destination-port
entropy. This distinction is important because explanation metadata improves
operator trust and post-event auditability, but it should not be confused with
the policy threshold itself.

% \begin{figure}[htbp]
% \centering
% \fbox{\parbox{0.9\linewidth}{
% \centering
% \textbf{SHAP audit flow in \padonap{}}\\[3pt]
% Window feature vector $\rightarrow$ XGBoost prediction
% $\rightarrow$ top-$k$ SHAP feature attributions
% $\rightarrow$ AI event payload
% $\rightarrow$ policy audit record.
% }}
% \caption{SHAP explanation workflow for auditable DDoS policy decisions. The
% explanation record connects the predicted class, confidence, top contributing
% features, and selected response tier.}
% \label{fig:shap-audit-flow}
% \end{figure}

In the current evaluation, SHAP is treated as audit metadata rather than as a
separately benchmarked runtime component. Measuring explanation-generation
latency and validating operator-facing explanation quality are left as future
work.

%----------------------------------------------------------------------
%----------------------------------------------------------------------

%----------------------------------------------------------------------
\section{Chapter Summary}
\label{sec:bg-summary}
%----------------------------------------------------------------------

This chapter introduced the background concepts required for \padonap{}. It
first described DDoS attacks in Data Center Networks (DCNs) and explained why
shared-resource infrastructure increases the operational impact of such
attacks. It then introduced CICDDoS2019, window-based feature extraction, and
the six-class traffic label space used by the implemented Track~A detector. The chapter also
reviewed SDN, NFV, CNFs, and Service Function Chaining as the networking
foundation for dynamic mitigation.

The chapter then described ONAP closed-loop automation and clarified the
boundary between \padonap{}'s lightweight implementation and a full ONAP
deployment. It further introduced the dual-track AI pipeline, consisting of
XGBoost-based immediate DDoS classification and CNN+BiLSTM+Attention
multi-horizon forecasting. Finally, it discussed SHAP explainability,
response-tier policy translation, LP-based SLA-aware allocation, and CNF
readiness delay.

These foundations show that proactive DDoS defense in DCNs is not only a
machine learning classification problem. It is also an orchestration problem
involving forecast lead time, policy decision making, CNF readiness, traffic
enforcement, resource feasibility, and evidence-based evaluation. The next
chapter builds on these foundations by reviewing related work on ML-based DDoS
detection, traffic forecasting, NFV/SFC mitigation, and closed-loop
orchestration, and by identifying the research gap addressed by this thesis.
%══════════════════════════════════════════════════════════════════════
%  CHAPTER 3 — Related Work
%══════════════════════════════════════════════════════════════════════
%======================================================================
\chapter{Related Work}
\label{ch:related}
%======================================================================

This chapter reviews prior work related to the design of \padonap{}. The
review focuses on four research areas: ML-based DDoS detection and forecasting,
SDN/NFV/CNF-based mitigation, closed-loop orchestration, and SLA-aware resource
management. Rather than comparing systems only by detection accuracy, this
chapter examines whether existing approaches provide forecasting lead time,
orchestration support, explainable policy translation, and resource-aware
mitigation. Based on this analysis, the chapter identifies the research gaps
addressed by this thesis.

%----------------------------------------------------------------------
\section{ML-Based DDoS Detection and Forecasting}
\label{sec:related-ml}
%----------------------------------------------------------------------

Machine-learning-based DDoS detection has been widely studied because ML models
can identify statistical traffic patterns that are difficult to capture using
static rules or manually configured thresholds. Recent surveys show that
statistical analysis, machine learning, deep learning, hybrid methods,
policy-based mitigation, and moving target defense are all active directions
in DDoS detection and mitigation research, especially in SDN environments
\cite{hill2024ddos_sdn_survey}. More recent AI-oriented surveys further show
that the field has moved from rule-based detection toward AI-based detection,
dataset-driven evaluation, adversarial robustness, and AI-assisted mitigation
\cite{apostu2025ai_ddos_survey}. However, the same survey also emphasizes
that many AI-DDoS studies still treat detection delay and mitigation as
emerging practical problems rather than as a fully integrated closed-loop
orchestration workflow.

A representative example is the ensemble online machine-learning approach of
Alashhab \textit{et al.}, which studies DDoS detection and mitigation in SDN
using an ensemble online ML model \cite{alashhab2024ensemble_oml}. The system
uses online learning to adapt to changing traffic conditions and is evaluated
with SDN-based experiments. This direction is important because online learning
can reduce dependence on static offline models. However, the main focus remains
current-state detection and SDN-level response. The work does not connect
forecasted attack risk to an NFV/CNF lifecycle decision, nor does it evaluate
whether the warning arrives early enough to hide orchestration delay.

Temporal models have also been explored for DDoS prediction and classification.
Kumar \textit{et al.} evaluate multivariate LSTM variants on CICDDoS2019 using
22 selected CICFlowMeter features \cite{kumar2025multivariate}. Hernandez
\textit{et al.} study multivariate time-series classification for high-speed
DDoS detection using LSTM, TCN, Transformer models, and SHAP-based feature
analysis \cite{hernandez2025realtime_timeseries}. Thaqi \textit{et al.}
combine online machine learning with threat intelligence for timely DDoS
prediction and report low-latency inference in a monitoring pipeline
\cite{thaqi2026online_threat_intel}. These works strengthen the case for
time-series or streaming analytics, but they remain primarily detection or
prediction systems. They do not evaluate a multi-tier orchestration policy that
uses several forecast horizons to pre-position mitigation functions before a
reactive trigger occurs.

Therefore, existing ML-based DDoS research provides strong foundations for
classification and temporal traffic analysis, but it does not fully solve the
operational problem addressed in this thesis: converting AI warnings into
timely and auditable mitigation actions under orchestration delay.

%----------------------------------------------------------------------
\section{SDN, NFV, and CNF-Based DDoS Mitigation}
\label{sec:related-sdn-nfv}
%----------------------------------------------------------------------

SDN and NFV have been widely used in DDoS mitigation research because they
provide programmability and deployment flexibility. SDN enables dynamic traffic
steering, filtering, and redirection, while NFV allows security functions such
as firewalls, rate limiters, intrusion detection systems, and scrubbers to be
deployed as software-based functions. In more recent systems, cloud-native
network functions extend this model by packaging network functions as
containers managed through Kubernetes or similar platforms.

SDN-based DDoS mitigation systems commonly use controller logic to detect
abnormal traffic and install forwarding or filtering rules. Wang
\textit{et al.} propose a two-stage SDN method that first detects suspicious
switch behavior using multi-dimensional traffic characteristics and then traces
and isolates attack sources \cite{wang2024multidimensional_sdn}. This
demonstrates the value of combining traffic features with programmable
mitigation. However, the response is still driven by observed abnormality and
threshold-style evidence, not by multi-horizon forecasts used to prepare
mitigation before the attack peak.

NFV-based and cloud-native systems extend this idea by dynamically deploying,
scaling, or composing mitigation functions. These approaches demonstrate that
programmable networks can support automated DDoS response. However, many of
them remain reactive by construction: the mitigation loop starts only after
attack traffic crosses a threshold, after a classifier detects an attack, or
after a controller observes abnormal switch behavior. This leaves an exposure
window that includes monitoring, detection, policy evaluation, orchestration,
VNF/CNF deployment, readiness checking, and traffic steering.

Elastic DDoS mitigation architectures such as CADA show that DDoS defense can
be treated as a flexible service composition and deployment problem. CADA
proposes a flexible and elastic DDoS mitigation architecture and uses
software-driven concepts to construct mitigation service chains
\cite{cada_2025}. This is relevant to \padonap{} because it recognizes that
DDoS mitigation is not only a detection problem, but also a service deployment
and orchestration problem. However, such systems do not explicitly combine
multi-horizon forecasting with ONAP-inspired policy translation, explanation
metadata, and SLA-aware allocation.

The key limitation of many SDN/NFV/CNF mitigation systems is therefore not
that they cannot mitigate attacks. Rather, the limitation is that they often
treat mitigation as a reactive action and lack a proactive preparation tier
that can be activated before the attack becomes severe. They also do not
explicitly account for the time required to make a mitigation function
enforceable. \padonap{} addresses this problem by using forecasted risk to
prepare mitigation before the reactive threshold alone would trigger
enforcement.

%----------------------------------------------------------------------
\section{ Orchestration and ONAP-Inspired Automation}
\label{sec:related-orchestration}
%----------------------------------------------------------------------

Network orchestration has matured significantly through ETSI NFV MANO and
related platforms. The ETSI NFV MANO framework defines the roles of the NFV
Orchestrator, VNF Manager, and Virtualized Infrastructure Manager for managing
the lifecycle of virtualized network services. ONAP is one of the most
comprehensive open-source platforms inspired by these principles, providing
support for service design, telemetry collection, policy evaluation,
orchestration, and closed-loop automation \cite{onap_docs}.

In ONAP-style closed-loop automation, telemetry and events from the managed
environment are collected, analyzed, evaluated by policy logic, and translated
into orchestration actions. This architecture is attractive for DDoS defense
because mitigation requires coordination between monitoring, analytics, policy,
deployment, and enforcement. However, deploying full ONAP is complex and often
unnecessary for small research testbeds. For this reason, many research systems
implement selected ONAP-like roles rather than the full production ONAP stack.

\padonap{} follows this lightweight interpretation. It adopts a DCAE-like
collector, a Policy-like decision component, an SO-like orchestration adapter,
and Kubernetes-based CNF enforcement. This design allows the thesis to study
the closed-loop control path without claiming full ONAP deployment. The
important research question is not whether a full ONAP platform is installed,
but whether AI outputs can be transformed into timely, explainable, and
resource-aware orchestration actions.

Compared with prior orchestration-oriented work, the distinguishing feature of
\padonap{} is the explicit connection between multi-horizon forecasting and
closed-loop mitigation. Existing prediction systems may generate early alerts,
and elastic mitigation systems may optimize deployment after a request is
formed. The missing link is the policy interface between the two: how a
forecasted risk at a specific horizon becomes a graduated orchestration action
with measurable lifecycle cost.

%----------------------------------------------------------------------
\section{SLA-Aware and Resource-Aware Mitigation}
\label{sec:related-sla}
%----------------------------------------------------------------------

DDoS mitigation in Data Center Networks must consider not only attack
suppression but also service quality for legitimate traffic. In shared-resource
environments, a mitigation action may protect one service while consuming
bandwidth, CPU, queueing capacity, or CNF processing resources that are also
needed by other services or tenants. Therefore, a practical mitigation system
should reason about resource feasibility and service priorities.

Many DDoS detection systems report accuracy, precision, recall, or F1-score,
but they do not evaluate whether mitigation decisions preserve service-level
requirements. Similarly, many SDN/NFV mitigation systems focus on traffic
redirection or filtering, but they do not explicitly model bandwidth floors or
allocation feasibility during active mitigation. This creates a gap between
attack detection and service-preserving defense.

\padonap{} addresses this issue through an LP-based SLA-aware allocation
mechanism. The allocator is not described as providing unconditional SLA
guarantees. Instead, it reasons about configured bandwidth floors and available
capacity. When the allocation problem is feasible, the system can preserve
configured floors; when it is infeasible, the policy layer can record the
condition and select an alternative response. This distinction is important
because it avoids overstating the system's capabilities while still providing
a formal mechanism for resource-aware mitigation.

%----------------------------------------------------------------------
%----------------------------------------------------------------------
\section{Research Gaps}
\label{sec:gaps}
%----------------------------------------------------------------------

Based on the reviewed literature, this thesis identifies four research gaps
that remain relevant to the current state of DDoS research.

\begin{enumerate}
    \item \textbf{G1 --- Prediction is rarely coupled to enforceable
    orchestration actions.}
    Recent LSTM and multivariate time-series studies show that temporal
    traffic modelling can improve DDoS prediction or real-time classification
    \cite{kumar2025multivariate,hernandez2025realtime_timeseries}. Online
    threat-intelligence systems also demonstrate low-latency prediction
    pipelines \cite{thaqi2026online_threat_intel}. However, these systems
    primarily produce detections, predictions, or alerts. They do not evaluate
    a mitigation orchestrator that converts multiple forecast horizons into
    concrete response tiers such as pre-positioning, mitigation, or isolation.
    \padonap{} addresses this gap by making multi-horizon forecasting an input
    to the policy engine rather than only an offline prediction result.

    \item \textbf{G2 --- Most mitigation mechanisms remain reactive and lack a
    proactive preparation stage.}
    SDN mitigation systems can trace, filter, or isolate attack traffic after
    abnormal behavior is observed \cite{wang2024multidimensional_sdn}. Elastic
    architectures such as CADA improve service-chain deployment and resource
    utilization for DDoS mitigation \cite{cada_2025}. These contributions are
    valuable, but the defense action is still mainly organized around response
    after detection or after a mitigation request has already been formed. What
    is missing is an explicit proactive stage that uses early warning evidence
    to prepare a lightweight mitigation function before the reactive trigger.
    \padonap{} addresses this gap by reporting proactive lead time and by
    comparing lightweight \tier{2} preparation against heavier \tier{3}
    mitigation.

    \item \textbf{G3 --- Explainability is mostly model-side, not
    policy-audit-side.}
    Recent detection work uses SHAP to identify discriminative traffic
    features and improve interpretability or feature selection
    \cite{hernandez2025realtime_timeseries}. This is useful, but it does not
    by itself explain why an orchestration tier was selected. In a closed-loop
    defense system, the audit record must connect model evidence, forecast
    horizon, confidence, resource state, and the selected mitigation action.
    \padonap{} addresses this gap by carrying explanation metadata as part of
    the AI output payload consumed by the policy layer.

    \item \textbf{G4 --- Resource-aware mitigation is not usually tied to
    tenant-level service floors.}
    Elastic mitigation work recognizes that DDoS defense is a deployment and
    resource-allocation problem, for example by optimizing service-chain
    placement or resource utilization \cite{cada_2025}. However, many
    detection and SDN mitigation studies still report accuracy, response time,
    or filtering success without checking whether the selected action preserves
    configured service floors for legitimate tenants. \padonap{} addresses this
    gap through an LP-based allocation check and reports SLA-floor preservation
    only under the evaluated capacity settings.

\end{enumerate}

%----------------------------------------------------------------------
\section{Chapter Summary}
\label{sec:related-summary}
%----------------------------------------------------------------------

This chapter reviewed representative work on ML-based DDoS detection,
temporal DDoS analytics, SDN/NFV/CNF mitigation, closed-loop orchestration,
and SLA-aware resource management. The review shows that existing systems are
strong in individual dimensions, such as detection accuracy, temporal traffic
analysis, or mitigation elasticity. However, they do not fully address the
combined operational problem targeted by this thesis: using multi-horizon AI
forecasting to support timely, explainable, and resource-aware mitigation
under orchestration delay.

These gaps motivate the design of \padonap{}, which integrates immediate
detection, multi-horizon forecasting, explainable policy translation,
closed-loop orchestration, and SLA-aware allocation. The next chapter presents
the proposed architecture and methodology in detail.
%══════════════════════════════════════════════════════════════════════
%  CHAPTER 4 — Proposed Method
%══════════════════════════════════════════════════════════════════════

%======================================================================
\chapter{Proposed Method}
\label{ch:proposed}
%======================================================================

This chapter presents the proposed \padonap{} framework. \padonap{}
stands for \textit{Proactive Attack Defense and Orchestration Model}.
The framework is designed to reduce the effective
mitigation delay of reactive DDoS defense in Data Center Networks (DCNs)
by combining near-real-time detection, multi-horizon forecasting,
explainable policy translation, SLA-aware resource allocation, and
Kubernetes-based CNF enforcement.

The chapter is organized as follows. Section~\ref{sec:approach} explains
the design rationale and the proactive mitigation principle.
Section~\ref{sec:arch-overview} presents the closed-loop architecture.
Section~\ref{sec:threat-model} defines the threat model and operational
assumptions. Section~\ref{sec:features} describes the AI components and
feature engineering. Section~\ref{sec:policy-sla} presents the policy
engine and SLA-aware allocation mechanism. Finally,
Section~\ref{sec:proposed-summary} summarizes the proposed method.

%----------------------------------------------------------------------
\section{Approach Direction}
\label{sec:approach}
%----------------------------------------------------------------------
The core design idea of ​​\padonap{} is that the delay in mitigating a DDoS attack is not just a detection issue. In an NFV/CNF-based environment, a reactive defense loop must complete several sequential steps before mitigation becomes effective: telemetry data collection, feature extraction, AI inference, policy evaluation, orchestration, deployment or activation of CNF, readiness check, and traffic execution. This time period is referred to in this thesis as the attack-to-execution latency.

A fully reactive system only begins this process after the attack has been detected with sufficient confidence. As a result, malicious traffic may continue to affect the protected service while the mitigation function is being prepared. \padonap{} addresses this limitation by using forecasting as an operational signal. Instead of simply waiting for a reaction threshold, the system uses short-term predictions to prepare mild mitigation actions before the attack reaches its most damaging stage.

The proposed design is based on three principles:

\begin{enumerate}

\item \textbf{Preparation based on predictions:}

Multi-horizon predictions are used to estimate near-future attack risk. When the predicted risk increases, \padonap{} can trigger preparatory mitigation actions, such as pre-positioning or scaling a lightweight CNF-based mitigation function. This reduces the effective exposure window when the prediction is accurate, while avoiding immediate disruptive enforcement when the evidence remains uncertain.

\item \textbf{Mitigation by level instead of binary response:}

\padonap{} does not treat every alert as a direct mitigation command.Instead, it maps detection and prediction outputs to hierarchical response levels,including normal monitoring, alert, pre-location, mitigation, and isolation. This allows the system to differentiate between early warning,confirmed attack, and severe attack states.

\item \textbf{Resource-based decision-making:}

Mitigation actions consume resources. Activating rate limits, filters or isolation rules can affect bandwidth, CPU, queue capacity, or traffic used by legitimate services. Therefore, \padonap{} includes an LP-based SLA-aware allocateer that checks whether the configured minimum bandwidth level can be maintained at the selected mitigation level. When the allocation is not feasible, the policy layer logs that situation and may choose an alternative response.
\end{enumerate}

%----------------------------------------------------------------------
\section{System Architecture Overview}
\label{sec:arch-overview}
%----------------------------------------------------------------------

\padonap{} is organized as a four-stage ONAP-inspired closed-loop pipeline.
The architecture is shown in Figure~\ref{fig:arch}.

\begin{figure}[htbp]
    \centering
    \includegraphics[width=\textwidth]{thesis/figs/pipeline.png}
    \caption{Overview of the \padonap{} four-stage ONAP-inspired closed-loop
    architecture.}
    \label{fig:arch}
\end{figure}

\textbf{Stage S1: Telemetry and feature extraction.}
Traffic telemetry is collected from the data center network through a
gNMI/gRPC telemetry path. In the Mininet implementation, a gNMI-compatible
adapter samples Open vSwitch counters every 500~ms and publishes normalized
records to Apache Kafka. Apache Flink aggregates the stream into 5~s sliding
windows with a 1~s slide. Each emitted window is converted into a feature
vector \(\mathbf{x}_t \in \mathbb{R}^{17}\), representing traffic rate,
entropy, protocol distribution, TCP flag behaviour, packet-size statistics,
and timing behaviour.

The adapter boundary is described in the evaluation setup; the architecture
requires a structured telemetry stream, not native OVS gNMI support.

\textbf{Stage S2: Dual-track AI inference.}
The AI layer has two complementary tracks. Track~A performs immediate
multi-class DDoS detection using an XGBoost classifier. Track~B performs
multi-horizon binary attack forecasting using a CNN+BiLSTM+Attention model
at horizons of 60, 300, and 600~seconds. Track~A estimates the current
traffic class, while Track~B estimates the probability of a DDoS-class event
at multiple future horizons.

\textbf{Stage S3: DCAE-like event collection and policy decision.}
The inference results are encoded as an AI output event and sent to a
DCAE-like collector. The policy engine consumes this event, evaluates the
current detection confidence, forecasted risk, selected horizon, recent tier
history, and resource feasibility. It then selects a response tier.

\textbf{Stage S4: SO-like orchestration and CNF enforcement.}
The selected response tier is translated into an orchestration action by an
SO-like adapter. The adapter interacts with Kubernetes or K3s to deploy,
pre-position, scale, or update mitigation CNFs. Depending on the implemented
testbed, enforcement may include rate limiting, filtering, traffic steering,
or isolation rules.

This architecture follows ONAP closed-loop concepts while remaining a
lightweight implementation. Production ONAP behavior is claimed only for the
selected ONAP OOM experiment reported in Chapter~\ref{ch:eval}.

%----------------------------------------------------------------------
\section{Threat Model and Operational Assumptions}
\label{sec:threat-model}
%----------------------------------------------------------------------

The adversary is assumed to control multiple traffic sources and may launch
volumetric floods, protocol-exhaustion attacks, application-layer floods, or
gradual-ramp attacks. The attacker aims to degrade the availability of a
protected service by consuming bandwidth, connection state, processing
capacity, or application-layer resources.

This thesis focuses on orchestration-aware mitigation rather than telemetry
security. Therefore, the following assumptions are made:

\begin{itemize}
    \item Telemetry sources are available to the monitoring pipeline;
    \item The control plane is trusted and not compromised by the attacker;
    \item The Kubernetes or K3s cluster can execute orchestration requests;
    \item Mitigation CNFs are available as deployable or pre-positionable
    components;
    \item The attacker can change traffic intensity or attack type, but cannot
    directly modify policy rules or AI model parameters.
\end{itemize}

The system is designed to satisfy three operational safety requirements.

\begin{itemize}
    \item \textbf{SLA-aware allocation safety.}
    For a selected mitigation tier \(T\), the bandwidth allocation should
    satisfy each configured service or tenant floor whenever the allocation
    problem is feasible:
    \[
        a_u(T) \geq f_u, \quad \forall u.
    \]

    \item \textbf{Critical-service floor preservation.}
    For critical services, the allocation mechanism should preserve the
    configured minimum floor whenever sufficient capacity remains after
    mitigation overhead is considered.

    \item \textbf{Bounded tier oscillation.}
    The policy engine should avoid rapid oscillation between mitigation tiers.
    A hysteresis guard holds the selected tier for a configurable number of
    windows unless stronger evidence requires escalation.
\end{itemize}

These requirements are operational constraints rather than unconditional
guarantees. If available capacity is insufficient, the LP allocator reports
infeasibility to the policy layer rather than silently violating configured
floors.

%----------------------------------------------------------------------
\section{AI Components and Feature Engineering}
\label{sec:features}
%----------------------------------------------------------------------

The AI layer contains two tracks operating on the same window-level feature
representation. Track~A performs immediate detection, while Track~B provides
forecasting support for proactive mitigation.

\subsection{Feature Engineering}
\label{subsec:feature-engineering}

Each traffic window is represented by a 17-dimensional feature vector:

\[
    \mathbf{x}_t =
    [x_{t,1}, x_{t,2}, \ldots, x_{t,17}] \in \mathbb{R}^{17}.
\]

The features are grouped into traffic-rate, entropy, protocol-distribution,
TCP-flag, packet-size, and timing categories. Shannon entropy is computed as:

\[
    H(X) = -\sum_i p_i \log_2 p_i,
\]

where \(p_i\) is the empirical probability of observing value \(i\) in the
corresponding traffic distribution.

\begin{table}[htbp]
\centering
\small
\caption{Window-level feature vector used by \padonap{}.}
\label{tab:features}
\begin{tabular}{@{}cllp{0.42\linewidth}@{}}
\toprule
\# & Feature & Group & Description \\
\midrule
1  & \texttt{pkt\_rate}          & Rate     & Packets per second \\
2  & \texttt{byte\_rate}         & Rate     & Bytes per second \\
3  & \texttt{new\_flows\_rate}   & Rate     & New flows per second \\
4  & \texttt{avg\_pkt\_size}     & Packet   & Average packet size \\
5  & \texttt{pkt\_size\_std}     & Packet   & Standard deviation of packet size \\
6  & \texttt{flow\_duration\_mean} & Timing  & Mean flow duration \\
7  & \texttt{inter\_arrival\_mean} & Timing  & Mean inter-arrival time \\
8  & \texttt{inter\_arrival\_std}  & Timing  & Inter-arrival time variation \\
9  & \texttt{src\_ip\_entropy}   & Entropy  & Entropy of source IPs \\
10 & \texttt{dst\_ip\_entropy}   & Entropy  & Entropy of destination IPs \\
11 & \texttt{src\_port\_entropy} & Entropy  & Entropy of source ports \\
12 & \texttt{dst\_port\_entropy} & Entropy  & Entropy of destination ports \\
13 & \texttt{proto\_dist\_tcp}   & Protocol & Fraction of TCP traffic \\
14 & \texttt{proto\_dist\_udp}   & Protocol & Fraction of UDP traffic \\
15 & \texttt{proto\_dist\_icmp}  & Protocol & Fraction of ICMP traffic \\
16 & \texttt{syn\_ratio}         & TCP Flag & SYN ratio among TCP packets \\
17 & \texttt{fin\_ratio}         & TCP Flag & FIN ratio among TCP packets \\
\bottomrule
\end{tabular}
\end{table}

The feature set is an implementation choice for this thesis rather than a
universal feature standard. If the framework is deployed in another
environment, the feature set may be adjusted to match the available telemetry
sources and target attack classes.

\subsection{XGBoost-Based Multi-Class Detection}
\label{subsec:track-a}

Detection uses an XGBoost classifier for immediate DDoS classification. Given a
feature vector \(\mathbf{x}_t\), the model outputs a probability distribution
over \(K\) traffic classes:

\[
    \hat{\mathbf{y}}_t =
    [P(y_t=1|\mathbf{x}_t), \ldots, P(y_t=K|\mathbf{x}_t)].
\]

The predicted class is:

\[
    \hat{y}_t = \arg\max_k P(y_t=k|\mathbf{x}_t),
\]

and the confidence score is:

\[
    c_t = \max_k P(y_t=k|\mathbf{x}_t).
\]

In this thesis, the implemented detector uses the six final classes produced
by the Track~A preprocessing pipeline after dropping the empty
\texttt{WebDDoS} window class:
\texttt{BENIGN}, \texttt{Amplification}, \texttt{Syn}, \texttt{UDP-lag},
\texttt{TFTP}, and \texttt{Portmap}. The policy engine uses the predicted
class and confidence score as part of the response-tier selection process.

Model hyperparameters, training split, and measured inference latency are
reported in the evaluation chapter where reproducible experiment artifacts are
available.

\subsection{CNN+BiLSTM+Attention Multi-Horizon Forecasting}
\label{subsec:track-b}

The forecaster estimates a binary attack indicator at several future horizons
from a sequence of recent per-second aggregated windows. Given an input
sequence

\[
    \mathbf{X}_t =
    (\mathbf{x}_{t-L+1}, \ldots, \mathbf{x}_{t}),
\]

with $L=600$ (a 10-minute lookback at 1~s aggregation), the forecaster outputs
a probability vector

\[
    \hat{\mathbf{y}}_{t}^{\text{fcst}} =
    \bigl[\, \hat{p}_{t+60},\ \hat{p}_{t+300},\ \hat{p}_{t+600}\,\bigr]
    \in [0,1]^{3},
\]

where $\hat{p}_{t+h} \approx P\!\bigl(y_{t+h}=1 \mid \mathbf{X}_t\bigr)$ and
$y_{t+h}\!=\!1$ iff at least one DDoS-class flow lands inside second
$t+h$. The lookback-to-horizon ratio of $2{\times}$ for the longest horizon
follows the M-competition rule-of-thumb for sequence forecasters
\cite{hyndman2018forecasting}, and the 10-minute history is comparable to the
effective context window inferred from the PARD-SSM kill-chain
example~\cite{hiremath2026pardssm}.

\paragraph{Structure.}
After comparing five sequence families (Section~\ref{subsec:forecaster-comparison}),
\textbf{CNN+BiLSTM+Attention} is selected as the Track~B model. The
architecture combines three complementary blocks:

\begin{enumerate}
    \item \textbf{1-D CNN front-end} (two Conv1d $\rightarrow$ BatchNorm
          $\rightarrow$ ReLU layers, kernel size 3, 64 channels) that extracts
          local temporal patterns such as packet-rate spikes and short burst
          signatures.
    \item \textbf{Two-layer Bidirectional LSTM} (hidden size 64) that captures
          longer-range temporal structure over the 600-second lookback.
          Bidirectionality is permissible because both passes only see the
          lookback window, not the future to be predicted.
    \item \textbf{Attention pooling}, $\mathbf{w}_t =
          \operatorname{softmax}(\mathbf{W}_a \mathbf{z}_t)$ followed by
          $\bar{\mathbf{z}} = \sum_t w_t\,\mathbf{z}_t$, which learns
          per-timestep weights and replaces a uniform last-step readout. A
          small MLP head produces the three-horizon logits.
\end{enumerate}

This hybrid follows the line of work that combines convolutional and
recurrent components with attention for DDoS prediction
\cite{alghazzawi2021cnnbilstm, alanazi2022ensembleddos, cnnbilstmattn2025},
and is consistent with the meta-learner architecture of HAMC-ID
\cite{hamcid2025nature}.

\begin{figure}[htbp]
    \centering
    \includegraphics[width=0.98\linewidth]{thesis/figs/ai/cnn_bilstm_attention_arch.png}
    \caption{CNN+BiLSTM+Attention forecaster used in Track~B. A two-layer 1-D
    convolutional front-end extracts local temporal patterns from the
    600-second lookback; a two-layer bidirectional LSTM models longer-range
    structure; an attention pooling layer learns per-timestep weights before
    a small MLP head outputs the three horizon probabilities.}
    \label{fig:cnn-bilstm-attention-arch}
\end{figure}

\paragraph{Training.}
Each sample is built by sliding a window of length $L=600$ over the
chronologically ordered per-second feature stream with stride 10. The model
is trained for 18 epochs with AdamW (learning rate $5\!\times\!10^{-4}$, weight
decay $10^{-4}$), cosine annealing, gradient-norm clipping at 0.5, and a
binary Focal Loss \cite{lin2017focal} with $\alpha=0.25$ and $\gamma=2$ to
counter the residual positive-class imbalance after dataset-level filtering.
Numerical safety is enforced via \texttt{nan\_to\_num} guards before scaling,
inside every batch, and before the AUC computation, because CICFlowMeter can
emit negative IAT and division-by-zero throughput values when flow durations
collapse to zero.

\paragraph{Policy integration.}
The policy engine does not directly enforce mitigation based on a single
forecast probability. Instead, it combines forecasted risk, forecast horizon,
detection confidence, recent tier history, and resource feasibility to select
a response tier. This reduces the chance that a single uncertain forecast
triggers unnecessary disruptive mitigation. For this reason, the evaluation
reports F1@$0.5$ alongside threshold-optimal F1: the deployed policy uses a
fixed threshold rather than a per-horizon sweep.

%----------------------------------------------------------------------
\section{Policy Engine and SLA-Aware Allocation}
\label{sec:policy-sla}
%----------------------------------------------------------------------

The policy engine translates AI outputs into operational response tiers. It is
designed to combine immediate detection, future risk estimation, hysteresis,
and resource feasibility.

\subsection{Graduated Response Tiers}
\label{subsec:tiers}

The mitigation action space is discretized into five tiers, as shown in
Table~\ref{tab:tiers}. The table describes the logical tier design; the
threshold values used for evaluation are reported with the experimental
scenario configuration rather than treated as universal constants.

\begin{table}[htbp]
\centering
\small
\caption{Graduated response tiers used by \padonap{}.}
\label{tab:tiers}
\begin{tabular}{@{}cllp{0.42\linewidth}@{}}
\toprule
Tier & Name & Main signal & Action \\
\midrule
$T_0$ & Normal &
No attack evidence &
Continue monitoring. \\

$T_1$ & Alert &
Weak detection or forecast evidence &
Log event, increase observation sensitivity, and prepare policy context. \\

$T_2$ & Pre-position &
Forecasted near-future risk &
Prepare or pre-scale lightweight mitigation CNF without applying disruptive
traffic enforcement. \\

$T_3$ & Mitigation &
High-confidence attack evidence &
Activate filtering, rate limiting, scrubbing, or traffic steering. \\

$T_4$ & Isolation &
Severe or persistent attack evidence &
Apply stronger containment such as blackholing or strict isolation. \\
\bottomrule
\end{tabular}
\end{table}

The key role of \(T_2\) is to support proactive mitigation. If forecast
evidence is strong enough, the system can prepare a lightweight mitigation
function before the attack reaches its peak. If the forecast is incorrect, the
cost is limited by the fact that \(T_2\) is a preparation tier rather than a
strong enforcement tier.

\subsection{Hysteresis Guard}
\label{subsec:hysteresis}

Rapid tier changes can cause VNF or CNF churn, unstable traffic steering, and
unnecessary resource overhead. To reduce oscillation, \padonap{} uses a
hysteresis guard. A selected tier is maintained for at least \(k\) consecutive
windows unless stronger evidence requires escalation. De-escalation is allowed
only after the supporting attack or forecast signal has disappeared for the
configured hold period.

This mechanism helps distinguish transient fluctuations from sustained attack
signals. It also bounds the duration of unnecessary proactive preparation when
a forecast is not confirmed by later observations.

\subsection{LP-Based SLA-Aware Allocation}
\label{subsec:lp}

At each tier transition, the SLA-aware allocator checks whether the selected
mitigation action is feasible under the available capacity. Let \(C\) denote
the link or service capacity, \(o(T)\) denote the overhead introduced by tier
\(T\), \(a_u\) denote the bandwidth allocated to service or tenant \(u\),
\(f_u\) denote its configured minimum floor, \(d_u\) denote its maximum demand,
and \(w_u\) denote its priority weight.

The allocator solves the following linear program:

\begin{equation}
\begin{aligned}
  \max_{\mathbf{a}} &\quad \sum_{u} w_u a_u \\[4pt]
  \text{s.t.} &\quad \sum_{u} a_u \leq C - o(T), \\
  &\quad f_u \leq a_u \leq d_u, \quad \forall u.
\end{aligned}
\label{eq:lp}
\end{equation}

If the problem is feasible, the resulting allocation preserves all configured
floors. If the problem is infeasible, the allocator reports this condition to
the policy engine. The policy engine may then select an alternative response,
reduce best-effort allocation, prioritize critical services, or escalate to an
operator-visible event.

This design is SLA-aware decision support rather than an unconditional SLA
guarantee. Bandwidth floors can be preserved only when the capacity constraints
allow a feasible allocation.

%----------------------------------------------------------------------
\section{AI Output Payload}
\label{sec:payload}
%----------------------------------------------------------------------

To decouple AI inference from policy and orchestration logic, \padonap{} uses
a canonical AI output payload. The payload records the current detection
result, forecast outputs, explanation metadata, and resource context. A
simplified schema is shown below.

\begin{verbatim}
{
  "timestamp": "...",
  "window_id": "...",
  "current_class": "...",
  "detection_confidence": 0.0,
  "forecast": [
    {"horizon": "+60s",  "attack_probability": 0.0},
    {"horizon": "+300s", "attack_probability": 0.0},
    {"horizon": "+600s", "attack_probability": 0.0}
  ],
  "top_features": ["...", "..."],
  "resource_state": {
    "available_capacity": 0.0,
    "active_tier": "T0"
  }
}
\end{verbatim}

This schema allows the AI layer, policy engine, and orchestration adapter to
evolve independently. It also supports auditability because each policy
decision can be traced back to the AI evidence and resource state that
triggered it.

%----------------------------------------------------------------------
\section{Chapter Summary}
\label{sec:proposed-summary}
%----------------------------------------------------------------------

This chapter presented the proposed \padonap{} framework. The method is based
on the observation that DDoS mitigation in NFV/CNF environments is constrained
by attack-to-enforcement latency, not only by detection accuracy. To reduce the
effective exposure window, \padonap{} combines immediate XGBoost-based
detection with CNN+BiLSTM+Attention multi-horizon forecasting and maps the
resulting AI evidence into graduated response tiers.

The chapter also described how telemetry, AI inference, policy selection,
orchestration, and CNF enforcement are connected in the closed loop. Finally,
it introduced an LP-based SLA-aware allocation mechanism that preserves
configured bandwidth floors when the allocation problem is feasible and exposes
infeasible cases to the policy layer.

The next chapter evaluates the proposed design in terms of AI performance,
forecasting behaviour, proactive lead time, closed-loop latency, CNF readiness,
mitigation effectiveness, resource overhead, and SLA-aware allocation.
%══════════════════════════════════════════════════════════════════════
%  CHAPTER 5 — Experiments and Evaluation
%══════════════════════════════════════════════════════════════════════
\chapter{Experiments and Evaluation}
\label{ch:eval}

This chapter evaluates \padonap{} across three complementary settings:
a controlled synthetic harness, a Mininet fat-tree emulation testbed, and
a selected real ONAP OOM integration test. The synthetic harness replays
pre-constructed traffic windows through the orchestration logic in order to
isolate the behaviour of the AI, policy, SLA allocation, and VNF lifecycle
components. The Mininet testbed exercises the same pipeline using a
gNMI/gRPC telemetry adapter, Apache Kafka event transport, and Apache Flink
window aggregation over Open vSwitch counters. Finally, the ONAP OOM test
examines the latency impact of using a live ONAP control plane for a selected
mitigation scenario.

Four representative scenarios are evaluated: S1 (normal baseline), S2
(sudden UDP flood), S3 (gradual SYN ramp), and S4 (proactive versus reactive
head-to-head). These scenarios are used to evaluate three central aspects of
the proposed design: detection and forecasting behaviour, proactive lead-time
advantage, and SLA-aware allocation under mitigation. The scenarios are not
intended to exhaustively cover all DDoS attack types, but they provide
controlled evidence for the main operational mechanisms of \padonap{}.

%----------------------------------------------------------------------
\section{Experimental Setup}
\label{sec:implementation}
%----------------------------------------------------------------------

\subsection{Synthetic Evaluation Environment}

\paragraph{Evaluation methodology.}
The orchestration frequency guard is disabled during scenario replay to allow
consecutive tier transitions within a single scenario run. Two evaluation
paths are compared: the proposed AI-driven path and a rule-based threshold
baseline. The baseline applies seven static conditions on the same
17-dimensional feature vector used by \padonap{}.

Both paths share the same downstream stack: VNF lifecycle client, SFC manager,
SLA allocator, and latency tracker. This design keeps the comparison focused
on the detection and forecasting stage rather than on differences in the
downstream mitigation implementation.

\paragraph{VNF boot-time calibration.}
For the synthetic evaluation, VNF boot times are modelled using fixed
instantiation delays: rate-limiter (\tier{2}) at 0.5~s, traffic scrubber
(\tier{3}) at 6~s, and blackhole action (\tier{4}) at 0.2~s. The ONAP Service
Orchestrator is represented by a stub implementation that follows the expected
lifecycle interaction pattern without invoking a live Kubernetes deployment.
Section~\ref{sec:onap-real} reports a complementary experiment using a live
ONAP OOM deployment to compare the stub approximation against a real control
plane.

\subsection{Mininet Testbed Environment}

The Mininet emulation network uses a fat-tree topology with $k=4$, comprising
4 core, 8 aggregation, and 8 edge (Top-of-Rack) Open vSwitch (OVS) switches
for a total of 20 switches and 16 hosts ($k^3/4 = 16$), following the
Al-Fares fat-tree design~\cite{alfares2008scalable}. Each pod contains
4 hosts, 2 edge switches, and 2 aggregation switches; every aggregation
switch connects to $k/2=2$ core switches, giving 4 equal-cost core paths per
host pair. The full layout and an example attacker$\rightarrow$victim path
are shown in Figure~\ref{fig:testbed-fat-tree-k4}.

\begin{figure}[H]
\centering
 \includegraphics[width=0.95\linewidth]{thesis/figs/fattree.png}
\caption{Fat-tree $k=4$ Mininet testbed used for DDoS validation: 4 core
switches (\texttt{c1}\,--\,\texttt{c4}), 8 aggregation switches
(\texttt{ap\_0}, \texttt{ap\_1} per pod), 8 edge switches
(\texttt{ep\_0}, \texttt{ep\_1} per pod), and 16 hosts
(\texttt{h0}\,--\,\texttt{h15}). The red trace shows the cross-pod attack
path used in all scenarios: attacker \texttt{h0} (pod~0) traverses
\texttt{e0\_0} $\rightarrow$ \texttt{a0\_0} $\rightarrow$ \texttt{c1}
$\rightarrow$ \texttt{a3\_0} $\rightarrow$ \texttt{e3\_1} to reach victim
\texttt{h15} (pod~3).}
\label{fig:testbed-fat-tree-k4}
\end{figure}

The fat-tree symbol convention used throughout the thesis is summarized in
Table~\ref{tab:testbed-fat-tree-legend}, and the run-time configuration in
Table~\ref{tab:mininet-setup}.

\begin{table}[H]
\centering
\small
\caption{Fat-tree symbol convention for the testbed.}
\label{tab:testbed-fat-tree-legend}
\begin{tabular}{@{}llp{0.62\linewidth}@{}}
\toprule
Symbol & Role & Quantity \\
\midrule
\texttt{c}        & Core switch              & 4 switches (\texttt{c1}--\texttt{c4}) at the top, connecting every pod \\
\texttt{a}        & Aggregation switch       & 2 per pod (\texttt{ap\_0}, \texttt{ap\_1}); connects edge $\leftrightarrow$ core \\
\texttt{e}        & Edge (ToR) switch        & 2 per pod (\texttt{ep\_0}, \texttt{ep\_1}); connects host $\leftrightarrow$ aggregation \\
\texttt{h}        & Host                     & 16 hosts (\texttt{h0}--\texttt{h15}), 4 hosts per pod \\
\bottomrule
\end{tabular}
\end{table}

\begin{table}[H]
\centering
\small
\caption{Mininet testbed configuration for real-traffic validation.}
\label{tab:mininet-setup}
\begin{tabular}{@{}ll@{}}
\toprule
Parameter & Value \\
\midrule
Topology         & Fat-tree $k=4$: 4 core $+$ 8 aggregation $+$ 8 edge OVS switches, 16 hosts \\
Host addressing  & \texttt{10.\textit{pod}.\textit{edge}.\textit{host+1}/8} (e.g. \texttt{h0}=10.0.0.1, \texttt{h15}=10.3.1.2) \\
Link capacity    & 100~Mbps (host--edge), 200~Mbps (aggregation/core) \\
Telemetry source & \texttt{softflowd} per-host NetFlow v5 export to a Kafka producer (real path); \\
                 & a gNMI-compatible adapter over OVS counters is provided as a synthetic alternative \\
Message bus      & Apache Kafka (KRaft, NodePort 30992) \\
Stream processor & Apache Flink, 5~s window with 1~s slide \\
Attack generator & \texttt{hping3} (UDP flood and SYN flood modes) \\
Background load  & \texttt{iperf3} TCP, approximately 95~Mbps \\
Feature output   & 17-dimensional feature vector per Flink window \\
AI models        & Same checkpoint as synthetic runs; no retraining \\
VNF lifecycle    & Stub-simulated for S1--S4; live ONAP OOM for the S2-ONAP integration run \\
\bottomrule
\end{tabular}
\end{table}

\paragraph{Attack path.}
Attack traffic is generated from host \texttt{h0} in pod~0 (IP
\texttt{10.0.0.1}) toward victim host \texttt{h15} in pod~3 (IP
\texttt{10.3.1.2}), traversing two pod boundaries through a core switch as
highlighted by the red trace in Figure~\ref{thesis/fig:testbed-fat-tree-k4}. This
choice maximizes the number of fat-tree hops the attack must cross
($\texttt{h0} \rightarrow \texttt{e0\_0} \rightarrow \texttt{a0\_0}
\rightarrow \texttt{c1} \rightarrow \texttt{a3\_0} \rightarrow \texttt{e3\_1}
\rightarrow \texttt{h15}$), ensuring that aggregation and core switches
participate in both the offered traffic measurement and any later
OVS-based enforcement action.

\paragraph{Local--remote pipeline.}
The full \padonap{} pipeline is split between a local Mininet host and a
remote Kubernetes cluster running ONAP OOM, as shown in
Figure~\ref{fig:testbed-remote-pipeline}. Only two flows cross the
local$\leftrightarrow$remote boundary: (i)~the raw telemetry stream
(\texttt{pad.telemetry.raw}) published to the remote Kafka broker over TCP,
and (ii)~an HTTP \texttt{/metrics} poll used by the local
\texttt{RemoteTierPoller} to sample tier-decision metadata from the remote
metrics service. The complete ONAP control loop (DCAE AI signals, Policy
tier decision, SO CNF lifecycle, CNF/VNF action) remains entirely inside
the remote cluster, so the local machine never participates in the
orchestration control path.

\begin{figure}[H]
\centering
\includegraphics[width=0.95\linewidth]{thesis/figs/local.png}
\caption{\padonap{} local--remote pipeline. Local Mininet captures
NetFlow~v5 from \texttt{softflowd} on every host, the Kafka-producing
collector forwards records to the remote Kafka broker, and the remote
Flink-driven AI orchestrator feeds DCAE/Policy/SO inside ONAP OOM. The
local \texttt{RemoteTierPoller} only samples post-decision metrics for
evaluation reports; it does not participate in the closed loop.}
\label{fig:testbed-remote-pipeline}
\end{figure}

\paragraph{Telemetry implementation.}
The real-traffic pipeline uses per-host \texttt{softflowd} agents that
export NetFlow~v5 records to a local Kafka producer, which then publishes
to the remote Kafka topic. Apache Flink consumes the Kafka stream on the
remote cluster, applies a 5~s sliding window with a 1~s slide, and emits
the 17-dimensional feature vector used for AI inference and tier decision.
The earlier gNMI-compatible adapter, which reads OVS port and flow
counters directly every 500~ms, is retained as an alternative telemetry
path for synthetic validation runs that do not require NetFlow
infrastructure. Both adapters serialize records with timestamp, switch
identifier, port identifier, packet count, byte count, protocol counters,
and TCP flag counters, so Flink can derive rates, entropy values, protocol
fractions, packet-size statistics, and timing features regardless of the
chosen capture mechanism. Thus, the statement ``gNMI-compatible'' means
that the adapter exposes a structured telemetry stream with gNMI/YANG-like
field semantics for the synthetic experiment; it does not mean that OVS
natively provides a production gNMI server.

Reported byte-rate values in the Mininet attack snapshots are interpreted as
offered or telemetry-observed aggregate attack load estimates before
mitigation, not as guaranteed delivered goodput on a single 100~Mbps host-edge
link. This distinction is necessary because the attack generator can offer
traffic above the configured bottleneck capacity, while the emulated links
still bound the traffic actually delivered to the victim.

%----------------------------------------------------------------------
\section{AI Model Performance}
\label{sec:metrics}
%----------------------------------------------------------------------

Table~\ref{tab:ai-metrics} summarizes the offline performance of both AI
components. Track~A (detection) is evaluated on CICDDoS2019 using a temporal
80/20 split with a 20-window purge gap; Track~B (forecasting) is evaluated on
CICIDS2017 (Sharafaldin et al., 2018a) using a chronological per-file 80/20
split. CICIDS2017 is the most-used DDoS forecasting benchmark in the
literature (29\% of surveyed studies in Mittal et al., 2022) and is the same
dataset used by the contemporary SOTA early-warning work PARD-SSM
\cite{hiremath2026pardssm}. CICDDoS2019 was previously considered for Track~B
but was rejected because its $>$99\% attack ratio leaves no
benign$\rightarrow$attack transitions for a forecaster to learn from; CSE-CIC-IDS2018
was rejected because a persistence baseline reaches $\text{F1}\approx 0.96$,
reducing the task to current-state detection.

\begin{table}[H]
\centering
\small
\caption{AI component metrics. Track~A is evaluated on CICDDoS2019, Track~B
on CICIDS2017.}
\label{tab:ai-metrics}
\begin{tabular}{@{}lll@{}}
\toprule
Track & Metric & Value \\
\midrule
\multirow{5}{*}{XGBoost 6-class detector (CICDDoS2019)}
  & Accuracy              & 0.9731 \\
  & Balanced accuracy     & 0.8339 \\
  & Macro-F1              & 0.8257 \\
  & Macro AUC (OvR)       & 0.9997 \\
  & Fit time              & 11.2~s \\
\midrule
\multirow{6}{*}{CNN+BiLSTM forecaster(CICIDS2017)}
  & AUC at $t+60$~s   & 0.9445 \\
  & AUC at $t+300$~s  & 0.9540 \\
  & AUC at $t+600$~s  & 0.9593 \\
  & PR-AUC at $t+60/300/600$~s & 0.7576 / 0.7878 / 0.7977 \\
  & F1@0.5 at $t+60/300/600$~s & 0.7801 / 0.8353 / 0.7635 \\
  & Inference P99 latency      & 7.7~ms \\
\bottomrule
\end{tabular}
\end{table}

\begin{figure}[htbp]
    \centering
    \includegraphics[width=0.95\linewidth]{thesis/figs/ai/forecaster_f1_thr05.png}
    \caption{F1 at the fixed deployment threshold of 0.5 (deployment-realistic
    operating point). XGBoost-flat is competitive at 60 s but the recurrent
    baselines (StackedLSTM, GRU, TemporalCNN, Transformer$+$LSTM) collapse at
    longer horizons because their probability outputs are not well calibrated;
    CNN+BiLSTM+Attention and XGBoost-flat are the only learned models that
    remain above the persistence reference at every horizon.}
    \label{fig:forecaster-f1-thr05}
\end{figure}

\paragraph{Lift over persistence.}
A persistence baseline that simply replays the current-second label
($\hat{y}_{t+h}=y_t$) achieves mean AUC~$0.8554$ over the same three horizons,
so CNN+BiLSTM+Attention delivers an absolute $\Delta$AUC of
$+0.0963$, $+0.0985$, and $+0.0966$ at $t{+}60$, $t{+}300$, and $t{+}600$~s,
respectively (mean lift $\approx +9.7\%$). This is a non-trivial improvement
on CICIDS2017 because attack episodes already last tens of minutes (e.g., DoS
Hulk $\approx 2$~h, DDoS LOIT $\approx 71$~min, DoS Slowloris $\approx 23$~min)
and a copy-current-state forecaster is therefore a strong baseline. By
comparison, an earlier Track~B run on CSE-CIC-IDS2018 collapses to a $+0.02$
AUC gap over persistence, providing empirical evidence of the task degeneracy that
motivated the move to CICIDS2017.

\begin{figure}[htbp]
    \centering
    \includegraphics[width=0.9\linewidth]{thesis/figs/ai/forecaster_auc_vs_horizon.png}
    \caption{Per-horizon ROC-AUC of every forecaster candidate on CICIDS2017.
    CNN+BiLSTM+Attention is the only sequence model whose AUC stays above
    XGBoost-flat across all three horizons and the only model whose AUC
    grows monotonically with the horizon length.}
    \label{fig:forecaster-auc-vs-horizon}
\end{figure}

\paragraph{Why AUC increases with horizon.}
Unlike conventional time-series forecasting, where uncertainty grows with the
horizon, Track~B AUC increases mildly from $0.9445$ at $t{+}60$~s to
$0.9593$ at $t{+}600$~s. The cause is the long-tailed distribution of attack
durations on CICIDS2017: most positive samples at a 600-second horizon fall
inside an ongoing attack episode, whose features are highly predictive,
while positive samples at a 60-second horizon are concentrated near attack
transitions, which are intrinsically harder to localize in time. The
same evidence-accumulation effect is reported by PARD-SSM, whose ``Early
Detection Margin'' peaks around eight minutes. This means Track~B is best
understood as attack-episode state forecasting rather than
exact-onset prediction; the strategic 600-second horizon is therefore
the most reliable signal for proactive $T_2$ pre-positioning, while the
60-second tactical horizon is primarily used for short-loop confirmation.

\subsection{Classifier Benchmark for Track A}
\label{subsec:classifier-comparison}

To validate Track~A, an additional classifier comparison is conducted using
the same feature vector and the same temporal 80/20 split. The purpose of this
comparison is to quantify the trade-off between aggregate accuracy, minority
class behaviour, and explainability support.

Table~\ref{tab:classifier-comparison} compares the six Track~A baselines:
XGBoost, LightGBM, Random Forest, ExtraTrees, MLP,
and a lightweight 1D-CNN. All models are evaluated on the same held-out temporal
test windows after the six-class remap.

\begin{table}[H]
\centering
\small
\caption{Classifier comparison for Track~A using the six-class
CICDDoS2019 detector configuration.}
\label{tab:classifier-comparison}
\begin{tabular}{@{}lrrrrl@{}}
\toprule
Model & Accuracy & Bal. Acc. & Macro-F1 & Macro AUC & Observation \\
\midrule
1D-CNN       & 0.9604 & 0.8396 & 0.8402 & 0.9963 & Highest macro-F1 in this comparison \\
MLP          & 0.9722 & 0.8342 & 0.8275 & 0.9992 & Strong tabular neural baseline \\
ExtraTrees   & 0.9735 & 0.8347 & 0.8271 & 0.9997 & Highest accuracy among tabular baselines \\
RandomForest & 0.9734 & 0.8344 & 0.8262 & 0.9988 & Similar to ExtraTrees, slower fit \\
XGBoost      & 0.9731 & 0.8339 & 0.8257 & 0.9997 & Supports SHAP analysis \\
LightGBM     & 0.9731 & 0.8340 & 0.8254 & 0.9998 & Comparable to XGBoost \\
\bottomrule
\end{tabular}
\end{table}

The comparison shows that the lightweight 1D-CNN obtains the highest macro-F1, but
all models struggle on the \texttt{UDP-lag} minority class. XGBoost is retained
as the detector for SHAP analysis and policy-level explanations:
its aggregate accuracy is close to the other tabular baselines, its one-vs-rest
AUC is high, and its tree structure provides direct compatibility with the SHAP
analysis stage used later in the pipeline.

\subsection{Forecaster Selection: Sequence Inductive Bias vs. Persistence}
\label{subsec:forecaster-comparison}

To justify the choice of \textbf{CNN+BiLSTM+Attention} for Track~B, an
analogous controlled benchmark is conducted on CICIDS2017 using the same 600-second
lookback, the same per-second feature aggregation, the same chronological per-file 80/20 split, and the same focal-loss training schedule. Five sequence models are evaluated against two sanity baselines: a persistence forecaster ($\hat{y}_{t+h}=y_t$) and a tabular XGBoost classifier trained on the flattened lookback (denoted ``XGBoost-flat''). Including persistence is critical because, as discussed in Section~\ref{sec:metrics}, it is the strongest non-learned baseline whenever attack episodes persist over the forecast horizon. XGBoost-flat is included as the strongest non-sequence learner: it sees exactly the same features but discards the temporal ordering, so any sequence model that fails to beat it is providing no additional inductive bias.

\begin{table}[H]
\centering
\small
\caption{Forecaster comparison on CICIDS2017 (mean ROC-AUC across
$h\in\{60,300,600\}$~s).
$\Delta$AUC is the mean lift over the persistence baseline.}
\label{tab:forecaster-comparison}
\begin{tabular}{@{}lrrrl@{}}
\toprule
Model & Mean AUC & $\Delta$AUC vs Pers. & Params  \\
\midrule
\textbf{CNN+BiLSTM+Attention}  & \textbf{0.9526} & \textbf{+0.0972} & 190.6~K  \\
XGBoost-flat                   & 0.9433 & +0.0879 & 10.8~K  \\
TemporalCNN                    & 0.9377 & +0.0823 & 41.3~K   \\
Transformer$+$LSTM             & 0.9294 & +0.0740 & 540.0~K \\
StackedLSTM                    & 0.9039 & +0.0485 & 348.4~K\\
Persistence baseline           & 0.8554 &      & 0 \\
GRU                            & 0.8034 & $-0.0520$ & 263.4~K  \\
\bottomrule
\end{tabular}
\end{table}

\begin{figure}[htbp]
    \centering
    \includegraphics[width=0.95\linewidth]{thesis/figs/ai/forecaster_mean_auc.png}
    \caption{Mean ROC-AUC of every forecaster candidate on CICIDS2017
    averaged over $h\in\{60,300,600\}$~s. CNN+BiLSTM+Attention is the only
    sequence model that exceeds the XGBoost-flat baseline; the persistence
    baseline (dotted line) is shown for reference.}
    \label{fig:forecaster-mean-auc}
\end{figure}

\begin{figure}[htbp]
    \centering
    \includegraphics[width=0.95\linewidth]{thesis/figs/ai/forecaster_delta_auc.png}
    \caption{ROC-AUC lift over the persistence baseline at each horizon.
    CNN+BiLSTM+Attention is the only model with a near-constant positive
    lift across all three horizons.}
    \label{fig:forecaster-delta-auc}
\end{figure}

Three observations follow directly from
Table~\ref{tab:forecaster-comparison} and
Figures~\ref{fig:forecaster-mean-auc}--\ref{fig:forecaster-delta-auc}. First, only CNN+BiLSTM+Attention beats the XGBoost-flat baseline, which is the strict requirement for justifying any sequence model on this task; the other recurrent or attention-only designs do not extract enough value from the temporal ordering to outweigh XGBoost-flat's tabular boosting. Second, CNN+BiLSTM+Attention achieves a mean $+9.7\%$ AUC lift over persistence and is the only model that delivers a near-constant positive lift across all three horizons ($+0.0963$ / $+0.0985$ / $+0.0966$), suggesting that the convolutional front-end and attention pooling capture features that genuinely carry over the 600-second horizon. Third, model capacity is not predictive of forecast quality: the 540~K-parameter Transformer$+$LSTM and 348~K-parameter StackedLSTM both fall behind the 190~K CNN+BiLSTM+Attention and even the 41~K TemporalCNN, consistent with the line of work that finds compact
convolutional or hybrid models more sample-efficient than pure attention backbones on similarly sized intrusion-detection benchmarks
\cite{alanazi2022ensembleddos, cnnbilstmattn2025}.

CNN+BiLSTM+Attention is therefore selected for Track~B. The inference cost remains modest (190.6~K parameters, $\approx 7.7$~ms P99 latency on a single T4 GPU) and is dominated by the convolutional front-end, so the resulting forecaster is compatible with the same closed-loop budget used by the Track~A XGBoost detector.

%----------------------------------------------------------------------
\section{Scenario, Baseline, and Lead-Time Results}
\label{sec:results}
%----------------------------------------------------------------------

The four evaluated scenarios are designed to stress distinct axes of the orchestration system: no-attack robustness (S1), sudden volumetric attack requiring reactive escalation (S2), gradual-ramp attack exploiting the proactive forecast path (S3), and a controlled within-run comparison of both paths under mixed traffic (S4). Each scenario window spans 5~s, and tier decisions are made per window.

For comparison, a rule-based threshold baseline is implemented on the same 17-dimensional feature vector. The baseline uses fixed expert-defined rules:packet-rate thresholds at 800, 3{,}000, and 10{,}000~pkt/s for \tier{1}, \tier{2}, and \tier{3}; SYN ratio $>0.60$ and UDP protocol ratio $>0.85$ for \tier{3}; ICMP protocol ratio $>0.70$ for \tier{2}; and source-IP entropy $<0.80$ with packet rate $>500$~pkt/s for \tier{2}. These thresholds are fixed before scenario replay and are not tuned on the reported outcomes. The baseline shares the same downstream VNF lifecycle, SFC, SLA allocation, and latency tracking logic as \padonap{}, keeping the comparison focused on the AI detection and forecasting stage.

\begin{table}[H]
\centering
\small
\caption{Scenario, threshold-baseline, and lead-time summary. Positive lead time indicates that \padonap{} acted earlier than the baseline. Latency values are P50 over VNF instantiation events only.}
\label{tab:scenario-baseline-leadtime}
\begin{tabular}{@{}p{1.4cm}p{2.7cm}p{1.6cm}p{1.8cm}p{1.8cm}p{2.1cm}p{3.4cm}@{}}
\toprule
ID & Scenario & \padonap{} max & First proact. & Baseline first & Lead vs. baseline & Main outcome \\
\midrule
S1 & Normal baseline & $T_0$ & --- & --- & --- &
No false escalation; SLA floor preserved. \\
S2 & Sudden UDP flood & $T_3$ & --- & $w{=}32$ & --- &
Abrupt onset is handled reactively; \padonap{} \tier{3} P50 is 6{,}006~ms versus 6{,}009~ms for the baseline. \\
S3 & Gradual SYN ramp & $T_2$ & $w{=}42$ & $w{=}55$ & $+\mathbf{65}$~s &
Forecaster pre-positions \tier{2} at 505~ms and avoids \tier{3} under the evaluated policy configuration. \\
S4 & Mixed proactive/reactive & $T_3$ & $w{=}31$ & $w{=}32$ & $+\mathbf{5}$~s &
First action is similar in onset but lighter: \tier{2} at 505~ms rather than baseline \tier{3} at 6{,}007~ms. \\
\bottomrule
\end{tabular}
\end{table}

The table shows that proactive advantage depends on attack shape. S3 provides the clearest gain because the gradual SYN ramp gives the
CNN+BiLSTM+Attention track enough temporal evidence to act 65~s before the threshold baseline. S4 provides a smaller baseline lead time, but within the same run the
proactive \tier{2} action at window~31 precedes the later reactive \tier{3} action at window~58 by 135~s. S2 has no proactive advantage because the UDP flood starts abruptly; in this case, \padonap{} behaves as a reactive detector and preserves downstream latency parity with the threshold baseline.

%----------------------------------------------------------------------
%----------------------------------------------------------------------
\section{Mininet End-to-End Validation}
\label{sec:mininet-e2e}
%----------------------------------------------------------------------

Scenarios S1--S4 are replayed on the Mininet fat-tree testbed to validate the gNMI/Kafka/Flink telemetry path and OVS SFC rule installation. Telemetry sampling, Flink aggregation, and OVS rule updates are executed on the emulated network; VNF boot remains stub-simulated except for the selected ONAP integration test in Section~\ref{sec:onap-real}.Table~\ref{tab:mininet-validation-summary} summarizes the Mininet evidence.

\begin{table}[H]
\centering
\small
\caption{Condensed Mininet validation results. VNF boot values remain
stub-simulated except for the ONAP case reported in Section~\ref{sec:onap-real}.}
\label{tab:mininet-validation-summary}
\begin{tabular}{@{}p{1.8cm}p{3.0cm}p{1.4cm}p{2.0cm}p{2.0cm}p{1.5cm}p{2.4cm}@{}}
\toprule
Scenario & Key signal & Max tier & First trigger & VNF boot & OVS SFC & Result \\
\midrule
S1 Normal & iperf3 background & $T_0$ & --- & --- & --- & No escalation \\
S2 UDP flood & conf $=0.924$ & $T_3$ & first attack window & 6{,}012~ms & 11~ms & Reactive mitigation \\
S3 SYN ramp & $P(t+30)=0.749$ & $T_2$ & 70~s & 508~ms & 9~ms & +50~s lead \\
S4 Mixed & $P(t+30)=0.762$; conf $=0.939$ & $T_3$ & 40~s / 170~s & 508 / 6{,}022~ms & 9 / 12~ms & 10.8$\times$ boot gap \\
\bottomrule
\end{tabular}
\end{table}

The UDP-flood run is used as the representative latency breakdown because it exercises telemetry collection, immediate detection, policy escalation, stub scrubber instantiation, and OVS rule installation in a single reactive path.

\begin{table}[H]
\centering
\small
\caption{End-to-end latency breakdown for S2 in Mininet with stub VNF boot.}
\label{tab:mininet-s2-latency}
\begin{tabular}{@{}lr@{}}
\toprule
Stage & Measured latency \\
\midrule
gNMI sampling window                  & 5{,}000~ms \\
Kafka/Flink aggregation and feature extraction & 51~ms \\
XGBoost inference                     & 2.2~ms \\
CNN+BiLSTM+Attention inference        & 7.7~ms \\
Policy decision ($T_0 \to T_3$)       & 1~ms \\
VNF instantiation (stub, scrubber)    & 6{,}012~ms \\
OVS SFC rule installation             & 11~ms \\
\midrule
\textbf{Total first attack window}    & \textbf{11{,}085~ms} \\
\bottomrule
\end{tabular}
\end{table}

The telemetry-observed offered attack estimate drops from 541.7~Mbps to 41.3~Mbps after mitigation, representing a 92.4\% reduction in observed attack load. Across the Mininet runs, the measured OVS SFC insertion time stays in the 9--12~ms range. The results follow the same trend as the synthetic evaluation: abrupt floods require reactive mitigation, while gradual or mixed attacks allow earlier \tier{2} preparation.

%----------------------------------------------------------------------
\section{Real ONAP Integration Test}
\label{sec:onap-real}
%----------------------------------------------------------------------

The ONAP experiment is limited to Scenario~S2 and is used to compare the stub orchestration approximation against a live ONAP OOM control plane. The ONAP cluster is deployed on a server with 64 CPU cores, 128~GB RAM, and a 200~GB SSD, with Service Orchestrator, Message Router, Policy PAP, Policy PDP.

The measured value in this experiment is the total scrubber VNF instantiation time after XGBoost detects the UDP flood: 22{,}110~ms. Per-component ONAP logs are incomplete, so Table~\ref{tab:onap-real-vs-stub} reports an estimated stage breakdown to indicate how the measured total is distributed.

\begin{table}[H]
\centering
\small
\caption{VNF instantiation latency: stub approximation versus real
ONAP OOM deployment (S2 UDP flood, scrubber VNF). Real ONAP stage values are estimated from the recorded 22{,}110~ms total; future experiments should replace them with per-component logs.}
\label{tab:onap-real-vs-stub}
\begin{tabular}{@{}lrr@{}}
\toprule
Stage & Stub & Real ONAP OOM \\
\midrule
Feature extraction and AI inference & 61~ms   & 61~ms \\
Policy dispatch      & 1~ms    & 251~ms (est.) \\
SO lifecycle request handling       & ---     & 2{,}310~ms (est.) \\
Kubernetes pod scheduling           & ---     & 1{,}840~ms (est.) \\
Container start + readiness         & ---     & 17{,}709~ms (est.) \\
VNF instantiation total             & 6{,}012~ms & \textbf{22{,}110~ms} \\
OVS SFC rule installation           & 11~ms   & 11~ms \\
\midrule
\textbf{Total pipeline latency}     & \textbf{6{,}085~ms} & \textbf{22{,}182~ms} \\
\bottomrule
\end{tabular}
\end{table}

The result shows that the synthetic stub underestimates lifecycle latency in a live orchestration environment. This strengthens the motivation for proactive
\tier{2} pre-positioning, while also showing that future ONAP experiments need synchronized Policy, SO, Kubernetes, and enforcement timestamps before a precise per-component latency claim can be made.

Overall, the synthetic and Mininet results show consistent policy behavior: \tier{2} preparation costs approximately 0.5~s in the stub setting, while \tier{3} scrubber instantiation costs approximately 6~s. The ONAP OOM run replaces this stub assumption with a measured 22.1~s scrubber instantiation time for S2, so it is reported as an integration boundary rather than as a fully decomposed end-to-end ONAP benchmark.
%----------------------------------------------------------------------
\section{Chapter Summary}
\label{sec:eval-summary}
%----------------------------------------------------------------------

This chapter evaluated \padonap{} through synthetic replay, Mininet emulation, and one selected ONAP OOM integration test. The main findings are summarized below.

\begin{enumerate}
  \item Track~A (detection) evaluates six classifiers on the six-class CICDDoS2019 setup (Section~\ref{subsec:classifier-comparison}). The lightweight 1D-CNN achieves the highest macro-F1 (0.8402), while XGBoost provides the detector used for SHAP explanations, with 0.9731 accuracy, 0.8257 macro-F1, and 0.9997 macro AUC.

  \item Track~B (forecasting) selects \textbf{CNN+BiLSTM+Attention} on the basis of a seven-model comparison (Section~\ref{subsec:forecaster-comparison}) on CICIDS2017: mean ROC-AUC 0.9526 across $h\!\in\!\{60,300,600\}$~s, $+9.7\%$ over the persistence baseline, $+0.93\%$ over the strongest non-sequence learner (XGBoost-flat at 0.9433), and 7.7~ms P99 inference latency with 190.6~K parameters --- the only sequence model that beats XGBoost-flat across all horizons.

  \item In the synthetic and Mininet settings, all four scenarios complete successfully under the evaluated configuration. The Mininet runs validate the gNMI-compatible telemetry path, Flink window aggregation,and OVS enforcement behavior, while VNF boot remains stub-simulated except for the ONAP S2 run.

  \item The forecasting path provides useful lead time when the attack has temporal buildup. In S3, \padonap{} pre-positions \tier{2} 65~s before the threshold baseline escalates. In S4, the first proactive action occurs 135~s before the within-run reactive \tier{3} transition.

  \item Proactive \tier{2} preparation is materially cheaper than\tier{3} scrubber instantiation in the stub setting: about 0.5~s for the rate-limiter path versus about 6~s for the scrubber path. The real ONAP OOM S2 run increases the scrubber instantiation time to 22.1~s, highlighting the importance of reporting stub and live-orchestration measurements separately.
\end{enumerate}

\chapter*{Conclusion and Future Works}
\addcontentsline{toc}{chapter}{Conclusion and Future Works}
\label{ch:conclusion}

\section*{Conclusion}

This thesis presented \padonap{}, a proactive attack defense and orchestration model for DDoS mitigation in data center networks. The central contribution is the integration of traffic analytics, short-horizon forecasting, policy-tier translation, and resource-aware orchestration into one closed-loop workflow.

The results show that treating DDoS defense as an orchestration problem changes
the evaluation focus. Detection quality remains important, but the decisive question is whether an AI signal arrives early enough to prepare an enforceable response and whether that response remains feasible under shared-resource constraints. Under the evaluated settings, \padonap{} demonstrates this behavior for gradual and mixed attack scenarios while clearly separating measured, stub-simulated, and estimated orchestration evidence.

\section*{Limitations}

These results demonstrate the feasibility of proactive orchestration under controlled assumptions. Nevertheless, several limitations remain.. First, the synthetic VNF boot times and most
Mininet VNF lifecycle results are stub-simulated; they are useful for comparing policy behavior under controlled assumptions but do not replace full CNF lifecycle measurements. \\
Second, the real ONAP OOM experiment records the total
VNF instantiation time, while the per-component latency breakdown is estimated and requires timestamped logs in a future implementation. \\
Third, SLA preservation is verified through an LP allocation model under configured capacity settings rather than through full multi-tenant application-level SLA traffic measurements.

Finally, the evaluation covers four representative scenarios rather than the full diversity of DDoS attack strategies. The forecasting model is most useful for attacks with temporal buildup and provides limited proactive advantage for abrupt floods.

\section*{Future Works}

Future work should first complete a fully measured CNF lifecycle experiment, including Kubernetes scheduling, image pull, container readiness, CNF packet processing, and service-chain update logs. The ONAP OOM test should also be extended with synchronized timestamps from DCAE, Policy, SO,
Kubernetes, and the enforcement layer so that the real control-plane latency can be decomposed without estimation.

Second, the evaluation should include richer traffic profiles, additional DDoS datasets, and longer Mininet runs with multiple attackers and multiple victims. This would test whether the 17-dimensional window-level feature set generalizes across different data center traffic distributions.

Third, future experiments should measure tenant-facing application performance under active mitigation. Such measurements should include goodput, latency, packet loss, SLA floor violations, and recovery time after mitigation is removed.

Finally, the forecasting track can be improved through probability
calibration, additional temporal features, and online adaptation. These improvements would help reduce false proactive actions while preserving useful lead time for orchestration.
\newpage
\addcontentsline{toc}{chapter}{Bibliography}
\bibliographystyle{ieeetr}
\bibliography{refs}

\end{document}
